% Cloudra 源程序清单（软著申请用），tectonic 直接编译
\PassOptionsToPackage{fontset=none}{ctex}
\documentclass[10pt,a4paper]{article}
\usepackage{ctex}
\usepackage[top=16mm,bottom=16mm,left=20mm,right=20mm]{geometry}
\usepackage{listings}
\usepackage{fancyhdr}
\usepackage{fontspec}
\setmonofont{Menlo}
\setCJKmonofont{Songti SC}
\lstset{
  basicstyle=\ttfamily\fontsize{6.5}{8.3}\selectfont,
  columns=fullflexible,
  keepspaces=true,
  breaklines=true,
  breakatwhitespace=false,
  breakindent=0pt,
  postbreak=,
  tabsize=4,
  frame=none,
  xleftmargin=0pt,
  xrightmargin=0pt,
  aboveskip=0pt,
  belowskip=0pt,
}
\pagestyle{fancy}
\fancyhf{}
\fancyhead[L]{\small\ttfamily Cloudra 0.1.0rc1 源程序}
\fancyhead[R]{\small\ttfamily 第 \thepage{} 页}
\renewcommand{\headrulewidth}{0.4pt}
\setlength{\headheight}{14pt}
\begin{document}
\setcounter{page}{1}
\begin{lstlisting}[firstnumber=1,numbers=left,numberstyle=\tiny\ttfamily,numbersep=5pt]
cloudra/__init__.py:1	"""Cloudra - AI Agent collaboration framework for chemical research."""
cloudra/__init__.py:2	
cloudra/__init__.py:3	from importlib import import_module
cloudra/__init__.py:4	
cloudra/__init__.py:5	__version__ = "0.1.0rc1"
cloudra/__init__.py:6	
cloudra/__init__.py:7	# 兼容垫片：包结构重排（core/services/persistence/base/infra）前的
cloudra/__init__.py:8	# `from cloudra import X` 老路径全部继续可用（惰性加载，不提前拉起重依赖）。
cloudra/__init__.py:9	# 执行栈（cloudra.core.agent / message_history）已删除，不再留入口。
cloudra/__init__.py:10	_LAZY_MODULES = {
cloudra/__init__.py:11	    # services
cloudra/__init__.py:12	    "session_store": "cloudra.services.session_store",
cloudra/__init__.py:13	    "nodes": "cloudra.services.nodes",
cloudra/__init__.py:14	    # persistence
cloudra/__init__.py:15	    "database": "cloudra.persistence.database",
cloudra/__init__.py:16	    "artifacts": "cloudra.persistence.artifacts",
cloudra/__init__.py:17	    "memory": "cloudra.persistence.memory",
cloudra/__init__.py:18	    "paths": "cloudra.persistence.paths",
cloudra/__init__.py:19	    "fs_utils": "cloudra.persistence.fs_utils",
cloudra/__init__.py:20	    "workspace": "cloudra.persistence.workspace",
cloudra/__init__.py:21	    # base
cloudra/__init__.py:22	    "session_class": "cloudra.base.session_class",
cloudra/__init__.py:23	    "skill_class": "cloudra.base.skill_class",
cloudra/__init__.py:24	    "tools_class": "cloudra.base.tools_class",
cloudra/__init__.py:25	    # infra
cloudra/__init__.py:26	    "config": "cloudra.infra.config",
cloudra/__init__.py:27	    "llm_client": "cloudra.infra.llm_client",
cloudra/__init__.py:28	    "streaming": "cloudra.infra.streaming",
cloudra/__init__.py:29	    "token_estimator": "cloudra.infra.token_estimator",
cloudra/__init__.py:30	    "web_search_native": "cloudra.infra.web_search_native",
cloudra/__init__.py:31	    "slurm": "cloudra.infra.slurm",
cloudra/__init__.py:32	    "compute_utils": "cloudra.infra.compute_utils",
cloudra/__init__.py:33	    "operations": "cloudra.infra.operations",
cloudra/__init__.py:34	    "yolo_guard": "cloudra.infra.yolo_guard",
cloudra/__init__.py:35	    "platform_utils": "cloudra.infra.platform_utils",
cloudra/__init__.py:36	    "runtime_profile": "cloudra.infra.runtime_profile",
cloudra/__init__.py:37	}
cloudra/__init__.py:38	
cloudra/__init__.py:39	__all__ = [*_LAZY_MODULES, "__version__"]
cloudra/__init__.py:40	
cloudra/__init__.py:41	
cloudra/__init__.py:42	def __getattr__(name: str):
cloudra/__init__.py:43	    """Keep legacy package module paths without importing heavy deps eagerly."""
cloudra/__init__.py:44	    module_path = _LAZY_MODULES.get(name)
cloudra/__init__.py:45	    if module_path is not None:
cloudra/__init__.py:46	        value = import_module(module_path)
cloudra/__init__.py:47	        globals()[name] = value
cloudra/__init__.py:48	        return value
cloudra/__init__.py:49	    raise AttributeError(f"module 'cloudra' has no attribute {name!r}")
cloudra/__init__.py:50	
cloudra/__init__.py:51	
cloudra/__init__.py:52	def __dir__() -> list[str]:
\end{lstlisting}
\clearpage
\begin{lstlisting}[firstnumber=53,numbers=left,numberstyle=\tiny\ttfamily,numbersep=5pt]
cloudra/__init__.py:53	    return sorted(set(globals()) | set(__all__))
cloudra/__main__.py:1	"""Allow running as `python -m cloudra`."""
cloudra/__main__.py:2	from cloudra.cli import main
cloudra/__main__.py:3	
cloudra/__main__.py:4	if __name__ == "__main__":
cloudra/__main__.py:5	    main()
cloudra/base/__init__.py:1	"""cloudra.base 子包。"""
cloudra/base/role_identity.py:1	"""DAG 存储词表与 v2 role_id 之间的唯一映射。
cloudra/base/role_identity.py:2	
cloudra/base/role_identity.py:3	DAG 变更日志、CallReview sidecar 与 process-node ``assignee`` 用的是历史 a-N
cloudra/base/role_identity.py:4	词表（``a-0`` = Coordinator，``a-1``..``a-8`` = 八个科研子 Agent，
cloudra/base/role_identity.py:5	``a-supervisor`` = 科学督导）；runtime_v2 的实时身份是语义 role_id
cloudra/base/role_identity.py:6	（``coordinator`` / ``small_mol`` …）。两套词表只在**存储/读取边界**互换，
cloudra/base/role_identity.py:7	映射表只此一张——之前 subagent_bridge 的 ``_AGENT_ID_TO_ROLE``、
cloudra/base/role_identity.py:8	``subagents/registry.py`` 的 a-N 列表与 launcher/swarm 提示词里的角色清单各写一份，
cloudra/base/role_identity.py:9	新增角色要同步多处。
cloudra/base/role_identity.py:10	
cloudra/base/role_identity.py:11	放置理由：runtime_v2 侧（RoleProfile）与工具侧、DAG 侧都要读它，
cloudra/base/role_identity.py:12	所以放在最底层的 ``cloudra.base``，不引入反向依赖。
cloudra/base/role_identity.py:13	"""
cloudra/base/role_identity.py:14	
cloudra/base/role_identity.py:15	from __future__ import annotations
cloudra/base/role_identity.py:16	
cloudra/base/role_identity.py:17	#: role_id → DAG 存储词表。新增角色只需在此登记一次 + 在 BUILTIN_ROLES 声明
cloudra/base/role_identity.py:18	#: （roles.py 里的 RoleProfile.legacy_agent_id 直接取自本表，不手抄）。
cloudra/base/role_identity.py:19	LEGACY_AGENT_ID_BY_ROLE: dict[str, str] = {
cloudra/base/role_identity.py:20	    "coordinator": "a-0",
cloudra/base/role_identity.py:21	    "explorer": "a-1",
cloudra/base/role_identity.py:22	    "writer": "a-2",
cloudra/base/role_identity.py:23	    "researcher": "a-3",
cloudra/base/role_identity.py:24	    "small_mol": "a-4",
cloudra/base/role_identity.py:25	    "md": "a-5",
cloudra/base/role_identity.py:26	    "dft": "a-6",
cloudra/base/role_identity.py:27	    "experiment": "a-7",
cloudra/base/role_identity.py:28	    "env": "a-8",
cloudra/base/role_identity.py:29	    "supervisor": "a-supervisor",
cloudra/base/role_identity.py:30	}
cloudra/base/role_identity.py:31	
cloudra/base/role_identity.py:32	#: 反向索引（一次构建）。
cloudra/base/role_identity.py:33	_ROLE_BY_LEGACY_AGENT_ID: dict[str, str] = {
cloudra/base/role_identity.py:34	    legacy: role for role, legacy in LEGACY_AGENT_ID_BY_ROLE.items()
cloudra/base/role_identity.py:35	}
cloudra/base/role_identity.py:36	
cloudra/base/role_identity.py:37	
cloudra/base/role_identity.py:38	def legacy_agent_id_for_role(role_id: str) -> str:
cloudra/base/role_identity.py:39	    """v2 role_id → DAG 词表 id。未知 role（或传入的已是 legacy id）原样透传。
cloudra/base/role_identity.py:40	
cloudra/base/role_identity.py:41	    透传是刻意的：历史库里存着 a-N，读取边界不能再把未知值改成别的身份。
cloudra/base/role_identity.py:42	    """
cloudra/base/role_identity.py:43	    return LEGACY_AGENT_ID_BY_ROLE.get(role_id, role_id)
cloudra/base/role_identity.py:44	
cloudra/base/role_identity.py:45	
\end{lstlisting}
\clearpage
\begin{lstlisting}[firstnumber=46,numbers=left,numberstyle=\tiny\ttfamily,numbersep=5pt]
cloudra/base/role_identity.py:46	def role_id_for_legacy_agent_id(legacy_agent_id: str) -> str:
cloudra/base/role_identity.py:47	    """DAG 词表 id → v2 role_id。未知值（或已是 role_id）原样透传。
cloudra/base/role_identity.py:48	
cloudra/base/role_identity.py:49	    仅用于把历史输入（旧提示词/旧工具参数里的 ``a-3``）归一到语义身份；新事件与
cloudra/base/role_identity.py:50	    新工具参数一律以 role_id 为准。
cloudra/base/role_identity.py:51	    """
cloudra/base/role_identity.py:52	    return _ROLE_BY_LEGACY_AGENT_ID.get(legacy_agent_id, legacy_agent_id)
cloudra/base/session_class.py:1	
cloudra/base/session_class.py:2	import os
cloudra/base/session_class.py:3	from datetime import datetime
cloudra/base/session_class.py:4	
cloudra/base/session_class.py:5	from cloudra.persistence.fs_utils import atomic_write_json
cloudra/base/session_class.py:6	
cloudra/base/session_class.py:7	
cloudra/base/session_class.py:8	class Session:
cloudra/base/session_class.py:9	    """内存态会话镜像；数据库（session 表）是唯一持久化事实来源。
cloudra/base/session_class.py:10	
cloudra/base/session_class.py:11	    默认不再写 per-session JSON 快照文件；调试时可设置环境变量
cloudra/base/session_class.py:12	    ``CLOUDRA_DEBUG_SESSION_JSON=1`` 恢复写文件行为（upload_json 仍可显式调用）。
cloudra/base/session_class.py:13	    """
cloudra/base/session_class.py:14	    id: str
cloudra/base/session_class.py:15	    path: str
cloudra/base/session_class.py:16	    create_time: datetime
cloudra/base/session_class.py:17	    state: str = 'sleeping'
cloudra/base/session_class.py:18	
cloudra/base/session_class.py:19	    def __init__(self, id: str, path: str) -> None:
cloudra/base/session_class.py:20	        self.create_time = datetime.now()
cloudra/base/session_class.py:21	        self.id = id
cloudra/base/session_class.py:22	        self.path = path
cloudra/base/session_class.py:23	        self.state = 'sleeping'
cloudra/base/session_class.py:24	        if os.environ.get("CLOUDRA_DEBUG_SESSION_JSON") == "1":
cloudra/base/session_class.py:25	            dir_name = os.path.dirname(self.path)
cloudra/base/session_class.py:26	            if dir_name:
cloudra/base/session_class.py:27	                os.makedirs(dir_name, exist_ok=True)
cloudra/base/session_class.py:28	            self.upload_json()
cloudra/base/session_class.py:29	
cloudra/base/session_class.py:30	
cloudra/base/session_class.py:31	    def to_json(self) -> dict:
cloudra/base/session_class.py:32	        create_time = self.create_time
cloudra/base/session_class.py:33	        if isinstance(create_time, datetime):
cloudra/base/session_class.py:34	            create_time = create_time.isoformat()
cloudra/base/session_class.py:35	        json_dict = {
cloudra/base/session_class.py:36	            'id': self.id,
cloudra/base/session_class.py:37	            'path': self.path,
cloudra/base/session_class.py:38	            'create_time': create_time,
cloudra/base/session_class.py:39	            'state': self.state,
cloudra/base/session_class.py:40	        }
cloudra/base/session_class.py:41	        return json_dict
cloudra/base/session_class.py:42	
cloudra/base/session_class.py:43	    def upload_json(self):
cloudra/base/session_class.py:44	        atomic_write_json(self.path, self.to_json())
cloudra/base/session_class.py:45	
\end{lstlisting}
\clearpage
\begin{lstlisting}[firstnumber=1,numbers=left,numberstyle=\tiny\ttfamily,numbersep=5pt]
cloudra/base/skill_class.py:1	import hashlib
cloudra/base/skill_class.py:2	import os
cloudra/base/skill_class.py:3	
cloudra/base/skill_class.py:4	import frontmatter as frt
cloudra/base/skill_class.py:5	
cloudra/base/skill_class.py:6	# 模块级缓存按 SKILL.md 内容 revision 失效。Phase 6 后 skill manifest
cloudra/base/skill_class.py:7	# 允许在进程运行期间更新，不能再假设目录在整个进程生命周期中静态不变。
cloudra/base/skill_class.py:8	_SKILL_DATAS_CACHE: dict[str, tuple[str, list[dict]]] = {}
cloudra/base/skill_class.py:9	
cloudra/base/skill_class.py:10	# NOTE: Skills 模型仅作为 skill frontmatter 的字段参照，当前未被实例化使用。
cloudra/base/skill_class.py:11	# 若后续需要校验 skill 元数据，可取消注释并补全默认值。
cloudra/base/skill_class.py:12	# class Skills(BaseModel):
cloudra/base/skill_class.py:13	#     name: str
cloudra/base/skill_class.py:14	#     description: str
cloudra/base/skill_class.py:15	#     path: str
cloudra/base/skill_class.py:16	#     scripts: list = []
cloudra/base/skill_class.py:17	#     references: list = []
cloudra/base/skill_class.py:18	
cloudra/base/skill_class.py:19	
cloudra/base/skill_class.py:20	class skill_loader:
cloudra/base/skill_class.py:21	
cloudra/base/skill_class.py:22	    # Skills are bundled with the package, not user data（base/ 上溯一级到包根）
cloudra/base/skill_class.py:23	    base_path: str = os.path.join(
cloudra/base/skill_class.py:24	        os.path.dirname(os.path.dirname(os.path.abspath(__file__))), 'skills'
cloudra/base/skill_class.py:25	    )
cloudra/base/skill_class.py:26	
cloudra/base/skill_class.py:27	    def __init__(self):
cloudra/base/skill_class.py:28	        # FIX: 原代码直接用了模块级变量 base_path（缺少 self.），在实例方法中访问类属性应加 self
cloudra/base/skill_class.py:29	        self.namelst = os.listdir(self.base_path)
cloudra/base/skill_class.py:30	
cloudra/base/skill_class.py:31	        self.pathlst = [os.path.join(self.base_path, skill_name) for skill_name in self.namelst]
cloudra/base/skill_class.py:32	
cloudra/base/skill_class.py:33	        self.md_pathlst = [os.path.join(path,'SKILL.md') for path in self.pathlst if os.path.isdir(path)]
cloudra/base/skill_class.py:34	
cloudra/base/skill_class.py:35	    def skill_datas_get(self) -> list[dict]:
cloudra/base/skill_class.py:36	        revision_hash = hashlib.sha256()
cloudra/base/skill_class.py:37	        for md in sorted(self.md_pathlst):
cloudra/base/skill_class.py:38	            revision_hash.update(md.encode("utf-8"))
cloudra/base/skill_class.py:39	            try:
cloudra/base/skill_class.py:40	                with open(md, "rb") as stream:
cloudra/base/skill_class.py:41	                    revision_hash.update(stream.read())
cloudra/base/skill_class.py:42	            except OSError:
cloudra/base/skill_class.py:43	                continue
cloudra/base/skill_class.py:44	        revision = revision_hash.hexdigest()
cloudra/base/skill_class.py:45	        cached = _SKILL_DATAS_CACHE.get(self.base_path)
cloudra/base/skill_class.py:46	        if cached is None or cached[0] != revision:
cloudra/base/skill_class.py:47	            mds = self.md_pathlst
cloudra/base/skill_class.py:48	            datas = []
cloudra/base/skill_class.py:49	            for md in mds:
cloudra/base/skill_class.py:50	                with open(md, 'r', encoding='utf-8') as f:
cloudra/base/skill_class.py:51	                    post = frt.load(f)
cloudra/base/skill_class.py:52	                    metadata = post.metadata
\end{lstlisting}
\clearpage
\begin{lstlisting}[firstnumber=53,numbers=left,numberstyle=\tiny\ttfamily,numbersep=5pt]
cloudra/base/skill_class.py:53	                    content = post.content
cloudra/base/skill_class.py:54	
cloudra/base/skill_class.py:55	                    # FIX: dict.update() 返回 None，原代码把 None 当数据 append 了
cloudra/base/skill_class.py:56	                    # 这里直接修改 metadata dict 后 append（frontmatter 的 metadata 是普通 dict）
cloudra/base/skill_class.py:57	                    metadata['content'] = content
cloudra/base/skill_class.py:58	                    datas.append(metadata)
cloudra/base/skill_class.py:59	
cloudra/base/skill_class.py:60	            _SKILL_DATAS_CACHE[self.base_path] = (revision, datas)
cloudra/base/skill_class.py:61	
cloudra/base/skill_class.py:62	        return _SKILL_DATAS_CACHE[self.base_path][1]
cloudra/base/skill_class.py:63	
cloudra/base/skill_class.py:64	    def list_skills(self) -> list[dict]:
cloudra/base/skill_class.py:65	        """Return list of all available skills with name and description."""
cloudra/base/skill_class.py:66	        datas = self.skill_datas_get()
cloudra/base/skill_class.py:67	        return [{"name": d.get("name", ""), "description": d.get("description", "")} for d in datas if d.get("name")]
cloudra/base/skill_class.py:68	
cloudra/base/skill_class.py:69	    def list_references(self, skill_name: str) -> list[str]:
cloudra/base/skill_class.py:70	        """Return relative paths of all available reference files for a given skill."""
cloudra/base/skill_class.py:71	        possible_dirs = ['references', 'reference', 'ref']
cloudra/base/skill_class.py:72	        refs = []
cloudra/base/skill_class.py:73	        for dir_name in possible_dirs:
cloudra/base/skill_class.py:74	            dir_path = os.path.join(self.base_path, skill_name, dir_name)
cloudra/base/skill_class.py:75	            if os.path.isdir(dir_path):
cloudra/base/skill_class.py:76	                for root, _, files in os.walk(dir_path):
cloudra/base/skill_class.py:77	                    for file in sorted(files):
cloudra/base/skill_class.py:78	                        if not file.startswith('.'):
cloudra/base/skill_class.py:79	                            rel = os.path.relpath(os.path.join(root, file), dir_path)
cloudra/base/skill_class.py:80	                            refs.append(rel)
cloudra/base/skill_class.py:81	        return sorted(refs)
cloudra/base/skill_class.py:82	
cloudra/base/skill_class.py:83	    def get_reference(self, skill_name, ref_name) -> str:
cloudra/base/skill_class.py:84	        # 支持多种可能的资料文件夹命名约定
cloudra/base/skill_class.py:85	        possible_dirs = ['references', 'reference', 'ref']
cloudra/base/skill_class.py:86	
cloudra/base/skill_class.py:87	        for dir_name in possible_dirs:
cloudra/base/skill_class.py:88	            ref_path = os.path.join(self.base_path, skill_name, dir_name, ref_name)
cloudra/base/skill_class.py:89	            if os.path.isfile(ref_path):
cloudra/base/skill_class.py:90	                with open(ref_path, 'r', encoding='utf-8') as f:
cloudra/base/skill_class.py:91	                    return f.read()
cloudra/base/skill_class.py:92	            elif os.path.isdir(ref_path):
cloudra/base/skill_class.py:93	                dir_filelst = sorted(os.listdir(ref_path))
cloudra/base/skill_class.py:94	                return f"Reference '{ref_name}' is a directory containing: {dir_filelst}. Please specify {ref_name}/<filename>."
cloudra/base/skill_class.py:95	
cloudra/base/skill_class.py:96	        return 'Reference not found.'
cloudra/base/skill_class.py:97	
cloudra/base/skill_class.py:98	    # FIX: 原拼写为 skill_dsatas_get（多打了个 s）
cloudra/base/skill_class.py:99	    def get_content(self, skill_name) -> str:
cloudra/base/skill_class.py:100	        datas = self.skill_datas_get()
cloudra/base/skill_class.py:101	        for data in datas:
cloudra/base/skill_class.py:102	            if data['name'] == skill_name:
cloudra/base/skill_class.py:103	                return data['content']
cloudra/base/skill_class.py:104	        return 'skill does not exist.'
\end{lstlisting}
\clearpage
\begin{lstlisting}[firstnumber=105,numbers=left,numberstyle=\tiny\ttfamily,numbersep=5pt]
cloudra/base/skill_class.py:105	
cloudra/base/skill_class.py:106	
cloudra/base/skill_class.py:107	
cloudra/base/skill_class.py:108	
cloudra/base/skill_class.py:109	
cloudra/base/skill_class.py:110	
cloudra/base/skill_class.py:111	
cloudra/base/skill_class.py:112	
cloudra/base/skill_class.py:113	
cloudra/base/skill_class.py:114	
cloudra/base/skill_class.py:115	
cloudra/base/skill_class.py:116	
cloudra/base/skill_class.py:117	
cloudra/base/skill_class.py:118	
cloudra/base/tool_actor.py:1	"""工具看到的调用者身份（ToolActorContext）。
cloudra/base/tool_actor.py:2	
cloudra/base/tool_actor.py:3	工具（bash/read/write/dag_*/goal/todo/…）只需要「谁在调我」的少量信息：工作
cloudra/base/tool_actor.py:4	目录、会话与项目键、权限位、编排模式。此前它们直接 import
cloudra/base/tool_actor.py:5	``cloudra.core.agent`` 的 ``Agent``——那个类带 LLM 调用能力，工具因此与已下线
cloudra/base/tool_actor.py:6	的执行栈耦合；字段被收窄成具体类还逼得执行器用 ``object.__setattr__`` 绕过
cloudra/base/tool_actor.py:7	pydantic 校验（见 ``runtime_v2/adapters.py``）。该执行栈（``Agent``/
cloudra/base/tool_actor.py:8	``Coordinator``/``Follower`` 与 ``core/message_history.py``）已随评审 5.4
cloudra/base/tool_actor.py:9	删除，这里只剩工具真正需要的表面。
cloudra/base/tool_actor.py:10	
cloudra/base/tool_actor.py:11	这里只声明**工具真正用到**的表面。生产实现是 ``runtime_v2.adapters`` 里包
cloudra/base/tool_actor.py:12	``AgentSession`` 的那个对象；测试可用任意鸭子类型对象。
cloudra/base/tool_actor.py:13	
cloudra/base/tool_actor.py:14	刻意不声明的字段：``agent_id`` 等 DAG 库里的历史 id（工具用
cloudra/base/tool_actor.py:15	``getattr(..., default)`` 兜底读取，属于迁移边界，不是新工具参数的首选身份）。
cloudra/base/tool_actor.py:16	"""
cloudra/base/tool_actor.py:17	
cloudra/base/tool_actor.py:18	from __future__ import annotations
cloudra/base/tool_actor.py:19	
cloudra/base/tool_actor.py:20	from typing import TYPE_CHECKING, Protocol, runtime_checkable
cloudra/base/tool_actor.py:21	
cloudra/base/tool_actor.py:22	if TYPE_CHECKING:
cloudra/base/tool_actor.py:23	    from cloudra.tools.core.permission_model import Permission
cloudra/base/tool_actor.py:24	
cloudra/base/tool_actor.py:25	
cloudra/base/tool_actor.py:26	@runtime_checkable
cloudra/base/tool_actor.py:27	class ToolActorContext(Protocol):
cloudra/base/tool_actor.py:28	    """调用方 Agent 暴露给工具的只读身份（不含任何执行能力）。"""
cloudra/base/tool_actor.py:29	
cloudra/base/tool_actor.py:30	    #: 会话 id（v2 为 ``as_*``；历史数据可能是旧的 a-N 组合 id）。
cloudra/base/tool_actor.py:31	    id: str
cloudra/base/tool_actor.py:32	    #: 调用者显示名；v2 下即 role_id（``coordinator`` / ``small_mol`` / …）。
cloudra/base/tool_actor.py:33	    name: str
cloudra/base/tool_actor.py:34	    #: 调用者类型：``coordinator`` / ``supervisor`` / ``follower``（DAG actor 判据）。
cloudra/base/tool_actor.py:35	    type: str
cloudra/base/tool_actor.py:36	    #: 会话键（v2 = project_id）：todos/goal 等 per-agent 存储的分区键。
cloudra/base/tool_actor.py:37	    conversation_id: str
cloudra/base/tool_actor.py:38	    #: 工具解析相对路径的根目录（绝不回落进程 cwd）。
\end{lstlisting}
\clearpage
\begin{lstlisting}[firstnumber=39,numbers=left,numberstyle=\tiny\ttfamily,numbersep=5pt]
cloudra/base/tool_actor.py:39	    working_directory: str
cloudra/base/tool_actor.py:40	    #: 编排模式字符串（standard / dag_record / dag_guided）。
cloudra/base/tool_actor.py:41	    orchestration_mode: str
cloudra/base/tool_actor.py:42	    #: 会话 id 别名；同角色并行实例（swarm）据此隔离各自的 todos/goal。
cloudra/base/tool_actor.py:43	    session_id: str
cloudra/base/tool_actor.py:44	
cloudra/base/tool_actor.py:45	    def permission_all(self) -> "Permission":
cloudra/base/tool_actor.py:46	        """调用者的权限位（IntFlag 投影，供 bash 等工具判定命令类别）。"""
cloudra/base/tool_actor.py:47	        ...
cloudra/base/tools_class.py:1	from contextlib import contextmanager
cloudra/base/tools_class.py:2	from contextvars import ContextVar
cloudra/base/tools_class.py:3	from typing import TYPE_CHECKING, Any, Awaitable, Callable, Optional
cloudra/base/tools_class.py:4	
cloudra/base/tools_class.py:5	from pydantic import BaseModel
cloudra/base/tools_class.py:6	from pydantic_core import PydanticUndefined
cloudra/base/tools_class.py:7	
cloudra/base/tools_class.py:8	from cloudra.base.session_class import Session
cloudra/base/tools_class.py:9	from cloudra.persistence.database import ConversationDB, get_database_path
cloudra/base/tools_class.py:10	
cloudra/base/tools_class.py:11	# HACK: 这里原本缺少 import，属于历史遗留问题；get_tools_lst 定义在 tools/__init__.py
cloudra/base/tools_class.py:12	from cloudra.tools import get_tools_lst
cloudra/base/tools_class.py:13	from cloudra.tools.core.permission_model import Permission
cloudra/base/tools_class.py:14	
cloudra/base/tools_class.py:15	if TYPE_CHECKING:
cloudra/base/tools_class.py:16	    from cloudra.infra.yolo_guard import YoloGuard
cloudra/base/tools_class.py:17	
cloudra/base/tools_class.py:18	# runtime_v2 的审批钩子：LegacyToolExecutor 在执行工具前总是设置；
cloudra/base/tools_class.py:19	# 未设置即生产外路径（v1/CLI/Flet 已删除），request_approval 直接抛错。
cloudra/base/tools_class.py:20	# 签名: async (**fields) -> tuple[bool, str]，字段同 request_approval 的关键字参数。
cloudra/base/tools_class.py:21	_APPROVAL_HANDLER: ContextVar[
cloudra/base/tools_class.py:22	    Optional[Callable[..., Awaitable[tuple[bool, str]]]]
cloudra/base/tools_class.py:23	] = ContextVar("cloudra_approval_handler", default=None)
cloudra/base/tools_class.py:24	
cloudra/base/tools_class.py:25	
cloudra/base/tools_class.py:26	@contextmanager
cloudra/base/tools_class.py:27	def use_approval_handler(handler: Callable[..., Awaitable[tuple[bool, str]]]):
cloudra/base/tools_class.py:28	    """Scope a v2 approval handler to one tool execution (per-task isolation)."""
cloudra/base/tools_class.py:29	    token = _APPROVAL_HANDLER.set(handler)
cloudra/base/tools_class.py:30	    try:
cloudra/base/tools_class.py:31	        yield
cloudra/base/tools_class.py:32	    finally:
cloudra/base/tools_class.py:33	        _APPROVAL_HANDLER.reset(token)
cloudra/base/tools_class.py:34	
cloudra/base/tools_class.py:35	
cloudra/base/tools_class.py:36	def _tool_db_path(subagent) -> str:
cloudra/base/tools_class.py:37	    """根据 tool 关联的 Agent 返回对应 conversation 的数据库路径。"""
cloudra/base/tools_class.py:38	    conv_id = getattr(subagent, 'conversation_id', 'default') if subagent else 'default'
cloudra/base/tools_class.py:39	    return get_database_path(conv_id)
cloudra/base/tools_class.py:40	
cloudra/base/tools_class.py:41	
cloudra/base/tools_class.py:42	def get_tools() -> list[dict]:
cloudra/base/tools_class.py:43	    """将工具类加载为 OpenAI function 格式的 dict 列表。"""
\end{lstlisting}
\clearpage
\begin{lstlisting}[firstnumber=44,numbers=left,numberstyle=\tiny\ttfamily,numbersep=5pt]
cloudra/base/tools_class.py:44	    entries = get_tools_lst()
cloudra/base/tools_class.py:45	    tools_list = []
cloudra/base/tools_class.py:46	    for entry in entries:
cloudra/base/tools_class.py:47	        tool_instance = entry.instantiate()
cloudra/base/tools_class.py:48	        tools_list.append(tool_instance.tool_schema_json())
cloudra/base/tools_class.py:49	    return tools_list
cloudra/base/tools_class.py:50	
cloudra/base/tools_class.py:51	
cloudra/base/tools_class.py:52	class input_schema(BaseModel):
cloudra/base/tools_class.py:53	    '''
cloudra/base/tools_class.py:54	    工具输入模式基类，定义工具需要的输入参数和格式。
cloudra/base/tools_class.py:55	
cloudra/base/tools_class.py:56	    子类用法示例：
cloudra/base/tools_class.py:57	        class MyInput(input_schema):
cloudra/base/tools_class.py:58	            query: str = Field(description="搜索关键词")
cloudra/base/tools_class.py:59	            max_results: int = Field(default=10, description="最大返回条数")
cloudra/base/tools_class.py:60	
cloudra/base/tools_class.py:61	    通过 to_parameters_schema() 可自动生成符合 OpenAI Function Calling 规范的 parameters 字典。
cloudra/base/tools_class.py:62	    '''
cloudra/base/tools_class.py:63	
cloudra/base/tools_class.py:64	    @classmethod
cloudra/base/tools_class.py:65	    def to_parameters_schema(cls) -> dict:
cloudra/base/tools_class.py:66	        '''
cloudra/base/tools_class.py:67	        生成符合 OpenAI Function Calling 的 parameters 字典。
cloudra/base/tools_class.py:68	        基于 Pydantic model_json_schema() 微调：
cloudra/base/tools_class.py:69	        - properties / required 已由 Pydantic 根据字段声明自动生成
cloudra/base/tools_class.py:70	        - 强制附加 additionalProperties: false，防止 LLM 传入未定义字段
cloudra/base/tools_class.py:71	        '''
cloudra/base/tools_class.py:72	        schema = cls.model_json_schema()
cloudra/base/tools_class.py:73	        parameters = {
cloudra/base/tools_class.py:74	            "type": "object",
cloudra/base/tools_class.py:75	            "properties": schema.get("properties", {}),
cloudra/base/tools_class.py:76	            "required": schema.get("required", []),
cloudra/base/tools_class.py:77	            "additionalProperties": False,
cloudra/base/tools_class.py:78	        }
cloudra/base/tools_class.py:79	        if "$defs" in schema:
cloudra/base/tools_class.py:80	            parameters["$defs"] = schema["$defs"]
cloudra/base/tools_class.py:81	        return parameters
cloudra/base/tools_class.py:82	
cloudra/base/tools_class.py:83	    @classmethod
cloudra/base/tools_class.py:84	    def from_args(cls, arg_dict: dict) -> "input_schema":
cloudra/base/tools_class.py:85	        '''从 LLM 返回的 arguments dict 构造输入实例，带 Pydantic 类型校验。'''
cloudra/base/tools_class.py:86	        return cls(**arg_dict)
cloudra/base/tools_class.py:87	
cloudra/base/tools_class.py:88	    def toClass(self, arg_dict: dict) -> "input_schema":
cloudra/base/tools_class.py:89	        '''
cloudra/base/tools_class.py:90	        兼容旧接口，供 reAct.py 调用。
cloudra/base/tools_class.py:91	        实际调用链通常为 tool.input_schema.toClass(tool_args)，此时 self 为 input_schema 子类本身。
cloudra/base/tools_class.py:92	        '''
cloudra/base/tools_class.py:93	        # 当通过类调用时（如 TextReader_input_schema.toClass(args)），self 即为类对象
cloudra/base/tools_class.py:94	        if isinstance(self, type) and issubclass(self, input_schema):
cloudra/base/tools_class.py:95	            return self(**arg_dict)
\end{lstlisting}
\clearpage
\begin{lstlisting}[firstnumber=96,numbers=left,numberstyle=\tiny\ttfamily,numbersep=5pt]
cloudra/base/tools_class.py:96	        return self.__class__(**arg_dict)
cloudra/base/tools_class.py:97	
cloudra/base/tools_class.py:98	    def toJson(self) -> dict:
cloudra/base/tools_class.py:99	        '''实例序列化为 dict，用于日志或上下文回传。'''
cloudra/base/tools_class.py:100	        return self.model_dump(mode="json")
cloudra/base/tools_class.py:101	
cloudra/base/tools_class.py:102	
cloudra/base/tools_class.py:103	class output_schema(BaseModel):
cloudra/base/tools_class.py:104	    '''
cloudra/base/tools_class.py:105	    工具输出模式基类，定义工具输出的格式。
cloudra/base/tools_class.py:106	    子类按普通 Pydantic Model 声明字段即可。
cloudra/base/tools_class.py:107	    '''
cloudra/base/tools_class.py:108	
cloudra/base/tools_class.py:109	    def toJson(self) -> dict:
cloudra/base/tools_class.py:110	        return self.model_dump(mode="json")
cloudra/base/tools_class.py:111	
cloudra/base/tools_class.py:112	
cloudra/base/tools_class.py:113	class Tool(BaseModel):
cloudra/base/tools_class.py:114	    '''
cloudra/base/tools_class.py:115	    工具基类。所有具体工具应继承此类。
cloudra/base/tools_class.py:116	
cloudra/base/tools_class.py:117	    关于 input_schema / output_schema：
cloudra/base/tools_class.py:118	    - 这里存放的是具体的 schema 类（如 TextReader_input_schema），而非实例。
cloudra/base/tools_class.py:119	    - 使用 Any 类型并开启 arbitrary_types_allowed，避免 Pydantic 对类对象做字段校验。
cloudra/base/tools_class.py:120	    - 子模块中覆盖时必须带类型注解，例如：
cloudra/base/tools_class.py:121	          input_schema: Any = TextReader_input_schema
cloudra/base/tools_class.py:122	    '''
cloudra/base/tools_class.py:123	    model_config = {"arbitrary_types_allowed": True}
cloudra/base/tools_class.py:124	
cloudra/base/tools_class.py:125	    # 唯一标识符
cloudra/base/tools_class.py:126	    tool_id: str = ''
cloudra/base/tools_class.py:127	
cloudra/base/tools_class.py:128	    # 工具名称（即 OpenAI function.name）
cloudra/base/tools_class.py:129	    tool_name: str = ''
cloudra/base/tools_class.py:130	
cloudra/base/tools_class.py:131	    # 工具描述。改造后应仅保留【功能描述】+【使用场景】，不再包含参数细节。
cloudra/base/tools_class.py:132	    # 参数说明通过 input_schema 的 Field(description="...") 提供，自动生成到 parameters 中。
cloudra/base/tools_class.py:133	    tool_description: str = ''
cloudra/base/tools_class.py:134	
cloudra/base/tools_class.py:135	    # 工具的输入模式类（存放类对象，非实例）
cloudra/base/tools_class.py:136	    input_schema: Any = input_schema
cloudra/base/tools_class.py:137	
cloudra/base/tools_class.py:138	    # 工具的输出模式类（存放类对象，非实例）
cloudra/base/tools_class.py:139	    output_schema: Any = output_schema
cloudra/base/tools_class.py:140	
cloudra/base/tools_class.py:141	    # 引用调用该工具的 Agent 实例
cloudra/base/tools_class.py:142	    subagent: Any = None
cloudra/base/tools_class.py:143	
cloudra/base/tools_class.py:144	    def is_read(self) -> bool:
cloudra/base/tools_class.py:145	        return False
cloudra/base/tools_class.py:146	
cloudra/base/tools_class.py:147	    def is_write(self) -> bool:
\end{lstlisting}
\clearpage
\begin{lstlisting}[firstnumber=148,numbers=left,numberstyle=\tiny\ttfamily,numbersep=5pt]
cloudra/base/tools_class.py:148	        return False
cloudra/base/tools_class.py:149	
cloudra/base/tools_class.py:150	    def is_manage(self) -> bool:
cloudra/base/tools_class.py:151	        return False
cloudra/base/tools_class.py:152	
cloudra/base/tools_class.py:153	    def is_script(self) -> bool:
cloudra/base/tools_class.py:154	        return False
cloudra/base/tools_class.py:155	
cloudra/base/tools_class.py:156	    def is_destructive(self) -> bool:
cloudra/base/tools_class.py:157	        return False
cloudra/base/tools_class.py:158	
cloudra/base/tools_class.py:159	    def tool_permission_all(self) -> Permission:
cloudra/base/tools_class.py:160	        """Return Permission flags for this tool."""
cloudra/base/tools_class.py:161	        perms = Permission.NONE
cloudra/base/tools_class.py:162	        if self.is_read():
cloudra/base/tools_class.py:163	            perms |= Permission.READ
cloudra/base/tools_class.py:164	        if self.is_write():
cloudra/base/tools_class.py:165	            perms |= Permission.WRITE
cloudra/base/tools_class.py:166	        if self.is_manage():
cloudra/base/tools_class.py:167	            perms |= Permission.MANAGE
cloudra/base/tools_class.py:168	        if self.is_script():
cloudra/base/tools_class.py:169	            perms |= Permission.SCRIPT
cloudra/base/tools_class.py:170	        if self.is_destructive():
cloudra/base/tools_class.py:171	            perms |= Permission.DESTRUCTIVE
cloudra/base/tools_class.py:172	        return perms
cloudra/base/tools_class.py:173	
cloudra/base/tools_class.py:174	    async def tool_call(self, user_input: input_schema, session: Session) -> output_schema:
cloudra/base/tools_class.py:175	        '''
cloudra/base/tools_class.py:176	        调用工具。子类必须重写此方法。
cloudra/base/tools_class.py:177	        执行结果以 output_schema 实例返回。
cloudra/base/tools_class.py:178	        '''
cloudra/base/tools_class.py:179	        raise NotImplementedError("子类必须实现 tool_call 方法")
cloudra/base/tools_class.py:180	
cloudra/base/tools_class.py:181	    async def safe_call(self, user_input: input_schema, session: Session) -> output_schema:
cloudra/base/tools_class.py:182	        '''
cloudra/base/tools_class.py:183	        带有通用异常捕获的工具调用包装器。
cloudra/base/tools_class.py:184	        工具执行失败时抛出包含工具名的 RuntimeError，由 reAct.act_result() 统一捕获并回传给 LLM。
cloudra/base/tools_class.py:185	        消息以 [Tool error] 前缀 + 工具名 + 恢复指引成形；react 层识别该前缀后不再二次套壳。
cloudra/base/tools_class.py:186	        '''
cloudra/base/tools_class.py:187	        try:
cloudra/base/tools_class.py:188	            return await self.tool_call(user_input, session)
cloudra/base/tools_class.py:189	        except Exception as e:
cloudra/base/tools_class.py:190	            error_msg = (
cloudra/base/tools_class.py:191	                f"[Tool error] {self.tool_name}: {str(e)}. "
cloudra/base/tools_class.py:192	                "Fix the arguments or use a different tool, then retry."
cloudra/base/tools_class.py:193	            )
cloudra/base/tools_class.py:194	            raise RuntimeError(error_msg) from e
cloudra/base/tools_class.py:195	
cloudra/base/tools_class.py:196	    async def ensure_safety(
cloudra/base/tools_class.py:197	        self,
cloudra/base/tools_class.py:198	        session: Session,
cloudra/base/tools_class.py:199	        details: str,
\end{lstlisting}
\clearpage
\begin{lstlisting}[firstnumber=200,numbers=left,numberstyle=\tiny\ttfamily,numbersep=5pt]
cloudra/base/tools_class.py:200	        tool_args: dict | None = None,
cloudra/base/tools_class.py:201	    ) -> tuple[bool, str]:
cloudra/base/tools_class.py:202	        '''
cloudra/base/tools_class.py:203	        安全检查：如果工具被标记为 destructive，则必须经过审批。
cloudra/base/tools_class.py:204	        默认通过数据库等待用户审批；若启用了 YOLO 模式，则调用独立 Agent 自动判断。
cloudra/base/tools_class.py:205	        返回 (是否通过, 拒绝原因)。通过时拒绝原因为空字符串。
cloudra/base/tools_class.py:206	        '''
cloudra/base/tools_class.py:207	        from cloudra.infra.yolo_guard import YoloGuard, is_yolo_enabled
cloudra/base/tools_class.py:208	
cloudra/base/tools_class.py:209	        risk_summary = next(
cloudra/base/tools_class.py:210	            (
cloudra/base/tools_class.py:211	                line.partition(":")[2].strip()
cloudra/base/tools_class.py:212	                for line in details.splitlines()
cloudra/base/tools_class.py:213	                if line.lower().startswith("action:") or line.startswith("【动作】")
cloudra/base/tools_class.py:214	            ),
cloudra/base/tools_class.py:215	            details.splitlines()[0] if details.splitlines() else "Dangerous tool call",
cloudra/base/tools_class.py:216	        )
cloudra/base/tools_class.py:217	
cloudra/base/tools_class.py:218	        if is_yolo_enabled():
cloudra/base/tools_class.py:219	            guard = YoloGuard()
cloudra/base/tools_class.py:220	            return await self._yolo_judge(session, tool_args or {}, guard, risk_summary=risk_summary)
cloudra/base/tools_class.py:221	
cloudra/base/tools_class.py:222	        db_path = _tool_db_path(self.subagent)
cloudra/base/tools_class.py:223	        candidate_diff = ""
cloudra/base/tools_class.py:224	        for marker in ("Unified Diff:\n", "【Unified Diff】:\n"):
cloudra/base/tools_class.py:225	            if marker in details:
cloudra/base/tools_class.py:226	                candidate_diff = details.partition(marker)[2]
cloudra/base/tools_class.py:227	                break
cloudra/base/tools_class.py:228	
cloudra/base/tools_class.py:229	        return await self.request_approval(
cloudra/base/tools_class.py:230	            db_path,
cloudra/base/tools_class.py:231	            session.id,
cloudra/base/tools_class.py:232	            details,
cloudra/base/tools_class.py:233	            tool_name=self.tool_name,
cloudra/base/tools_class.py:234	            normalized_arguments=tool_args or {},
cloudra/base/tools_class.py:235	            risk_summary=risk_summary,
cloudra/base/tools_class.py:236	            candidate_diff=candidate_diff,
cloudra/base/tools_class.py:237	        )
cloudra/base/tools_class.py:238	
cloudra/base/tools_class.py:239	    async def record_session_event(
cloudra/base/tools_class.py:240	        self,
cloudra/base/tools_class.py:241	        session: Session,
cloudra/base/tools_class.py:242	        event_type: str,
cloudra/base/tools_class.py:243	        data: dict,
cloudra/base/tools_class.py:244	    ) -> None:
cloudra/base/tools_class.py:245	        """尽力记录工具事件；遥测失败不能反向破坏工具主流程。"""
cloudra/base/tools_class.py:246	        from cloudra.services import session_store as _ss
cloudra/base/tools_class.py:247	
cloudra/base/tools_class.py:248	        try:
cloudra/base/tools_class.py:249	            db = ConversationDB.from_path(_tool_db_path(self.subagent))
cloudra/base/tools_class.py:250	            await _ss.arecord_event(db, session.id, event_type, data)
cloudra/base/tools_class.py:251	        except Exception:
\end{lstlisting}
\clearpage
\begin{lstlisting}[firstnumber=252,numbers=left,numberstyle=\tiny\ttfamily,numbersep=5pt]
cloudra/base/tools_class.py:252	            pass
cloudra/base/tools_class.py:253	
cloudra/base/tools_class.py:254	    async def _yolo_judge(
cloudra/base/tools_class.py:255	        self,
cloudra/base/tools_class.py:256	        session: Session,
cloudra/base/tools_class.py:257	        tool_args: dict,
cloudra/base/tools_class.py:258	        guard: 'YoloGuard',
cloudra/base/tools_class.py:259	        risk_summary: str = "",
cloudra/base/tools_class.py:260	    ) -> tuple[bool, str]:
cloudra/base/tools_class.py:261	        """调用 YoloGuard 做自动判断。"""
cloudra/base/tools_class.py:262	        approved, reason = await guard.judge(
cloudra/base/tools_class.py:263	            self.tool_name, tool_args, self.subagent, session, risk_summary=risk_summary
cloudra/base/tools_class.py:264	        )
cloudra/base/tools_class.py:265	        if approved:
cloudra/base/tools_class.py:266	            print("\033[92m[安全审计/YOLO]\033[0m 自动审批通过。")
cloudra/base/tools_class.py:267	            return True, ""
cloudra/base/tools_class.py:268	        else:
cloudra/base/tools_class.py:269	            reason_text = reason or "YOLO Agent 认为该操作存在风险"
cloudra/base/tools_class.py:270	            print(f"\033[91m[安全审计/YOLO]\033[0m 自动审批拒绝。原因: {reason_text}")
cloudra/base/tools_class.py:271	            return False, reason_text
cloudra/base/tools_class.py:272	
cloudra/base/tools_class.py:273	    async def request_approval(
cloudra/base/tools_class.py:274	        self,
cloudra/base/tools_class.py:275	        db_path: str,
cloudra/base/tools_class.py:276	        session_id: str,
cloudra/base/tools_class.py:277	        details: str,
cloudra/base/tools_class.py:278	        *,
cloudra/base/tools_class.py:279	        tool_name: str = "",
cloudra/base/tools_class.py:280	        normalized_arguments: dict | None = None,
cloudra/base/tools_class.py:281	        risk_summary: str = "",
cloudra/base/tools_class.py:282	        candidate_diff: str = "",
cloudra/base/tools_class.py:283	    ) -> tuple[bool, str]:
cloudra/base/tools_class.py:284	        handler = _APPROVAL_HANDLER.get()
cloudra/base/tools_class.py:285	        if handler is None:
cloudra/base/tools_class.py:286	            # v1 时代"SQLite + 事件总线等待"的审批回退已随旧服务层删除。
cloudra/base/tools_class.py:287	            raise RuntimeError(
cloudra/base/tools_class.py:288	                "no approval handler installed: runtime_v2 LegacyToolExecutor always sets one"
cloudra/base/tools_class.py:289	            )
cloudra/base/tools_class.py:290	        # runtime_v2 上下文：审批请求与等待由内核 broker 承担，
cloudra/base/tools_class.py:291	        # 不碰 conversation SQLite。
cloudra/base/tools_class.py:292	        return await handler(
cloudra/base/tools_class.py:293	            details=details,
cloudra/base/tools_class.py:294	            tool_name=tool_name,
cloudra/base/tools_class.py:295	            normalized_arguments=normalized_arguments or {},
cloudra/base/tools_class.py:296	            risk_summary=risk_summary,
cloudra/base/tools_class.py:297	            candidate_diff=candidate_diff,
cloudra/base/tools_class.py:298	        )
cloudra/base/tools_class.py:299	
cloudra/base/tools_class.py:300	    @classmethod
cloudra/base/tools_class.py:301	    def _get_field_default(cls, field_name: str):
cloudra/base/tools_class.py:302	        '''获取类字段的默认值，兼容 Pydantic v2（类级别无法直接访问普通字段）。'''
cloudra/base/tools_class.py:303	        field = cls.model_fields.get(field_name)
\end{lstlisting}
\clearpage
\begin{lstlisting}[firstnumber=304,numbers=left,numberstyle=\tiny\ttfamily,numbersep=5pt]
cloudra/base/tools_class.py:304	        if field and field.default is not PydanticUndefined:
cloudra/base/tools_class.py:305	            return field.default
cloudra/base/tools_class.py:306	        return None
cloudra/base/tools_class.py:307	
cloudra/base/tools_class.py:308	    def tool_schema_json(self) -> dict:
cloudra/base/tools_class.py:309	        '''
cloudra/base/tools_class.py:310	        返回标准 OpenAI function 格式，可直接用于 client.chat.completions.create(..., tools=[...])。
cloudra/base/tools_class.py:311	        优先读取实例级别的 tool_description，使动态生成的描述能够生效。
cloudra/base/tools_class.py:312	        '''
cloudra/base/tools_class.py:313	        # HACK: 这里用 getattr + or fallback 链是因为 Pydantic v2 的类字段默认值在实例上可能
cloudra/base/tools_class.py:314	        # 被 shadow，但某些情况下实例属性又还没覆盖。动态 description（如 Bash 的权限列表、
cloudra/base/tools_class.py:315	        # GetSkill 的 skills 列表）必须在 __init__ 之后才能确定，因此必须是实例方法而非类方法。
cloudra/base/tools_class.py:316	        # 如果后续有更好方式（如 descriptor / property），可替换掉这个 ugly fallback。
cloudra/base/tools_class.py:317	        description = getattr(self, 'tool_description', '') or self._get_field_default('tool_description') or ''
cloudra/base/tools_class.py:318	        return {
cloudra/base/tools_class.py:319	            "type": "function",
cloudra/base/tools_class.py:320	            "function": {
cloudra/base/tools_class.py:321	                "name": getattr(self, 'tool_name', '') or self._get_field_default('tool_name') or '',
cloudra/base/tools_class.py:322	                "description": description,
cloudra/base/tools_class.py:323	                "parameters": (self._get_field_default('input_schema') or input_schema).to_parameters_schema(),
cloudra/base/tools_class.py:324	            }
cloudra/base/tools_class.py:325	        }
cloudra/base/tools_class.py:326	
cloudra/base/tools_class.py:327	    def toJson(self) -> dict:
cloudra/base/tools_class.py:328	        '''
cloudra/base/tools_class.py:329	        实例级序列化，主要用于日志或调试。
cloudra/base/tools_class.py:330	        若需获取 OpenAI schema，请使用类方法 tool_schema_json()。
cloudra/base/tools_class.py:331	        '''
cloudra/base/tools_class.py:332	        return {
cloudra/base/tools_class.py:333	            "tool_id": self.tool_id,
cloudra/base/tools_class.py:334	            "tool_name": self.tool_name,
cloudra/base/tools_class.py:335	            "tool_description": self.tool_description,
cloudra/base/tools_class.py:336	        }
cloudra/base/tools_class.py:337	
cloudra/base/tools_class.py:338	    @classmethod
cloudra/base/tools_class.py:339	    def build_input(cls, arg_dict: dict) -> input_schema:
cloudra/base/tools_class.py:340	        '''
cloudra/base/tools_class.py:341	        将 LLM 返回的参数字典转换为 input_schema 实例。
cloudra/base/tools_class.py:342	        供 reAct.py 调用，替代原有的 tool.input_schema.toClass(tool_args)。
cloudra/base/tools_class.py:343	        '''
cloudra/base/tools_class.py:344	        schema_cls = cls._get_field_default('input_schema') or input_schema
cloudra/base/tools_class.py:345	        return schema_cls.from_args(arg_dict)
cloudra/benchmark/__init__.py:1	"""Benchmark module: benchmark task schema, task generators, and the scoring pipeline.
cloudra/benchmark/__init__.py:2	
cloudra/benchmark/__init__.py:3	Public surface:
cloudra/benchmark/__init__.py:4	- `schema`: TaskSpec / GoldSpec / RubricSpec / mechanism-graph models.
cloudra/benchmark/__init__.py:5	- `gen_easy_tasks`: generator turning ideapool/reaction_candidates/easy.md
cloudra/benchmark/__init__.py:6	  into ideapool/tasks/easy/*.json.
cloudra/benchmark/__init__.py:7	- `graph_similarity`: chemical-parameterized GED between two mechanism graphs.
cloudra/benchmark/__init__.py:8	- `integrity`: sandbox forensic scan + evidence binding (Layer 1).
cloudra/benchmark/__init__.py:9	- `selfconsistency`: physical self-consistency checks (Layer 2).
cloudra/benchmark/__init__.py:10	- `accuracy`: Gaussian-tolerance barrier scoring (Layer 3).
\end{lstlisting}
\clearpage
\begin{lstlisting}[firstnumber=11,numbers=left,numberstyle=\tiny\ttfamily,numbersep=5pt]
cloudra/benchmark/__init__.py:11	- `method_policy`: method-whitelist interception.
cloudra/benchmark/__init__.py:12	- `runner`: independent benchmark runner with its own data directory.
cloudra/benchmark/__init__.py:13	"""
cloudra/benchmark/__init__.py:14	
cloudra/benchmark/__init__.py:15	from cloudra.benchmark.schema import (  # noqa: F401
cloudra/benchmark/__init__.py:16	    AccuracyRubric,
cloudra/benchmark/__init__.py:17	    BondChangeSet,
cloudra/benchmark/__init__.py:18	    GoldMechanismGraph,
cloudra/benchmark/__init__.py:19	    GoldSpec,
cloudra/benchmark/__init__.py:20	    GraphSimilarityRubric,
cloudra/benchmark/__init__.py:21	    IntegrityRubric,
cloudra/benchmark/__init__.py:22	    MechanismEdge,
cloudra/benchmark/__init__.py:23	    MechanismNode,
cloudra/benchmark/__init__.py:24	    MethodPolicy,
cloudra/benchmark/__init__.py:25	    NodeRole,
cloudra/benchmark/__init__.py:26	    PromptSpec,
cloudra/benchmark/__init__.py:27	    RubricSpec,
cloudra/benchmark/__init__.py:28	    SelfConsistencyRubric,
cloudra/benchmark/__init__.py:29	    SystemSpec,
cloudra/benchmark/__init__.py:30	    TaskSpec,
cloudra/benchmark/__init__.py:31	    Tier,
cloudra/benchmark/__init__.py:32	)
cloudra/benchmark/__main__.py:1	"""python -m cloudra.benchmark — run the benchmark task bank.
cloudra/benchmark/__main__.py:2	
cloudra/benchmark/__main__.py:3	Default is a dry-run (no-op backend: sandboxes stay empty, every task scores
cloudra/benchmark/__main__.py:4	0 through the real scoring pipeline — useful for checking the harness itself).
cloudra/benchmark/__main__.py:5	`--backend runtime_v2` runs each task on the real in-process runtime_v2 kernel
cloudra/benchmark/__main__.py:6	with the method whitelist enforced.
cloudra/benchmark/__main__.py:7	
cloudra/benchmark/__main__.py:8	Example:
cloudra/benchmark/__main__.py:9	    python -m cloudra.benchmark --backend runtime_v2 --data-dir bench_out --only C1,D3
cloudra/benchmark/__main__.py:10	    python -m cloudra.benchmark --tasks-dir ideapool/tasks/easy --tasks-dir ideapool/tasks/impossible
cloudra/benchmark/__main__.py:11	"""
cloudra/benchmark/__main__.py:12	
cloudra/benchmark/__main__.py:13	from __future__ import annotations
cloudra/benchmark/__main__.py:14	
cloudra/benchmark/__main__.py:15	import argparse
cloudra/benchmark/__main__.py:16	import os
cloudra/benchmark/__main__.py:17	import sys
cloudra/benchmark/__main__.py:18	from pathlib import Path
cloudra/benchmark/__main__.py:19	
cloudra/benchmark/__main__.py:20	
cloudra/benchmark/__main__.py:21	def main(argv: list[str] | None = None) -> int:
cloudra/benchmark/__main__.py:22	    from cloudra.benchmark.runner import BenchmarkRunner, load_tasks
cloudra/benchmark/__main__.py:23	
cloudra/benchmark/__main__.py:24	    ap = argparse.ArgumentParser(prog="cloudra.benchmark")
cloudra/benchmark/__main__.py:25	    ap.add_argument("--tasks-dir", type=Path, action="append", default=None,
cloudra/benchmark/__main__.py:26	                    help="task JSON dir; repeatable (default: ideapool/tasks/easy)")
cloudra/benchmark/__main__.py:27	    ap.add_argument("--data-dir", type=Path, default=Path("benchmark_data"),
cloudra/benchmark/__main__.py:28	                    help="benchmark data dir (sandboxes + benchmark.db)")
cloudra/benchmark/__main__.py:29	    ap.add_argument("--backend", choices=["none", "runtime_v2"], default="none")
cloudra/benchmark/__main__.py:30	    ap.add_argument("--only", default="", help="comma-separated task ids (default: all)")
\end{lstlisting}
\clearpage
\begin{lstlisting}[firstnumber=31,numbers=left,numberstyle=\tiny\ttfamily,numbersep=5pt]
cloudra/benchmark/__main__.py:31	    ap.add_argument("--timeout", type=float, default=7200.0, help="per-task run timeout, seconds")
cloudra/benchmark/__main__.py:32	    ap.add_argument("--concurrency", type=int, default=1,
cloudra/benchmark/__main__.py:33	                    help="parallel tasks (threads, per-task event loop + sandbox; default 1 = serial)")
cloudra/benchmark/__main__.py:34	    ap.add_argument("--api-pool", type=Path, default=None,
cloudra/benchmark/__main__.py:35	                    help='JSON endpoint pool for anti-429: [{"base_url","api_key","label"}, ...]; '
cloudra/benchmark/__main__.py:36	                         "each LLM call picks a random endpoint, 429 fails over to the next")
cloudra/benchmark/__main__.py:37	    args = ap.parse_args(argv)
cloudra/benchmark/__main__.py:38	
cloudra/benchmark/__main__.py:39	    tasks = []
cloudra/benchmark/__main__.py:40	    for d in (args.tasks_dir or [Path("ideapool/tasks/easy")]):
cloudra/benchmark/__main__.py:41	        tasks.extend(load_tasks(d))
cloudra/benchmark/__main__.py:42	    if args.only:
cloudra/benchmark/__main__.py:43	        wanted = {t.strip() for t in args.only.split(",") if t.strip()}
cloudra/benchmark/__main__.py:44	        tasks = [t for t in tasks if t.task_id in wanted]
cloudra/benchmark/__main__.py:45	    if not tasks:
cloudra/benchmark/__main__.py:46	        print("no tasks matched", file=sys.stderr)
cloudra/benchmark/__main__.py:47	        return 2
cloudra/benchmark/__main__.py:48	
cloudra/benchmark/__main__.py:49	    backend = None
cloudra/benchmark/__main__.py:50	    if args.backend == "runtime_v2":
cloudra/benchmark/__main__.py:51	        # 运行时数据隔离必须在首次读取全局配置前完成；config_file 独立于
cloudra/benchmark/__main__.py:52	        # data_dir，LLM 凭据不受影响。YAML 里的 data_dir 键会覆盖环境变量，
cloudra/benchmark/__main__.py:53	        # 所以两处都钉死。
cloudra/benchmark/__main__.py:54	        isolated = (args.data_dir / "runtime").resolve()
cloudra/benchmark/__main__.py:55	        os.environ["CLOUDRA_DATA_DIR"] = str(isolated)
cloudra/benchmark/__main__.py:56	        from cloudra.infra.config import get_config
cloudra/benchmark/__main__.py:57	
cloudra/benchmark/__main__.py:58	        get_config().data_dir = isolated
cloudra/benchmark/__main__.py:59	        llm = None
cloudra/benchmark/__main__.py:60	        if args.api_pool is not None:
cloudra/benchmark/__main__.py:61	            from cloudra.benchmark.api_pool import PooledLLMProvider, load_endpoints
cloudra/benchmark/__main__.py:62	
cloudra/benchmark/__main__.py:63	            llm = PooledLLMProvider(load_endpoints(args.api_pool))
cloudra/benchmark/__main__.py:64	        from cloudra.benchmark.backend import RuntimeV2Backend
cloudra/benchmark/__main__.py:65	
cloudra/benchmark/__main__.py:66	        backend = RuntimeV2Backend(
cloudra/benchmark/__main__.py:67	            run_timeout_s=args.timeout,
cloudra/benchmark/__main__.py:68	            llm=llm,
cloudra/benchmark/__main__.py:69	            store_path=args.data_dir / "runtime" / "runtime-v2.db",
cloudra/benchmark/__main__.py:70	        )
cloudra/benchmark/__main__.py:71	
cloudra/benchmark/__main__.py:72	    # CLI 每次调用也是一个批次：batches 行 + 溯源快照（模型/模式/插件）+
cloudra/benchmark/__main__.py:73	    # runs 行带 batch_id——CLI 与 Web 结果在同一 schema 下可并排比较。
cloudra/benchmark/__main__.py:74	    from cloudra.benchmark.runner import provenance_config
cloudra/benchmark/__main__.py:75	    from cloudra.runtime_v2.schemas import new_id
cloudra/benchmark/__main__.py:76	
cloudra/benchmark/__main__.py:77	    batch_id = new_id("batch")
cloudra/benchmark/__main__.py:78	    runner = BenchmarkRunner(args.data_dir, backend=backend, batch_id=batch_id)
cloudra/benchmark/__main__.py:79	    runner.db.record_batch(batch_id, "running", {
cloudra/benchmark/__main__.py:80	        **provenance_config("cli", tasks),
cloudra/benchmark/__main__.py:81	        "tasks_dirs": [str(d) for d in (args.tasks_dir or [Path("ideapool/tasks/easy")])],
cloudra/benchmark/__main__.py:82	        "task_ids": [t.task_id for t in tasks],
\end{lstlisting}
\clearpage
\begin{lstlisting}[firstnumber=83,numbers=left,numberstyle=\tiny\ttfamily,numbersep=5pt]
cloudra/benchmark/__main__.py:83	        "concurrency": max(1, args.concurrency),
cloudra/benchmark/__main__.py:84	        "run_timeout_s": args.timeout,
cloudra/benchmark/__main__.py:85	    })
cloudra/benchmark/__main__.py:86	    status = "done"
cloudra/benchmark/__main__.py:87	    try:
cloudra/benchmark/__main__.py:88	        results = runner.run_all(tasks, concurrency=max(1, args.concurrency))
cloudra/benchmark/__main__.py:89	    except BaseException:
cloudra/benchmark/__main__.py:90	        status = "error"
cloudra/benchmark/__main__.py:91	        raise
cloudra/benchmark/__main__.py:92	    finally:
cloudra/benchmark/__main__.py:93	        runner.db.finish_batch(batch_id, status)
cloudra/benchmark/__main__.py:94	    for r in results:
cloudra/benchmark/__main__.py:95	        print(
cloudra/benchmark/__main__.py:96	            f"{r.task_id:>6}  final={r.final_score:.3f}  graph={r.graph_score:.3f} "
cloudra/benchmark/__main__.py:97	            f"selfcons={int(r.selfconsistency_passed)} accuracy={r.accuracy_score:.3f}"
cloudra/benchmark/__main__.py:98	        )
cloudra/benchmark/__main__.py:99	    return 0
cloudra/benchmark/__main__.py:100	
cloudra/benchmark/__main__.py:101	
cloudra/benchmark/__main__.py:102	if __name__ == "__main__":
cloudra/benchmark/__main__.py:103	    sys.exit(main())
cloudra/benchmark/accuracy.py:1	"""Layer 3 — accuracy: 以交付物为准的逐边能垒/性质评分（不读任何 LLM 散文）。
cloudra/benchmark/accuracy.py:2	
cloudra/benchmark/accuracy.py:3	口径（与题库 query 协议一致）：gold 机理图**逐边**承载参考答案能垒；评分边
cloudra/benchmark/accuracy.py:4	集合是**全部 gold 边**，逐边规则：
cloudra/benchmark/accuracy.py:5	
cloudra/benchmark/accuracy.py:6	- 带参考值（``barrier_kcal_mol`` 非 None）→ 双来源核对：要求两条独立来源
cloudra/benchmark/accuracy.py:7	  **都存在且一致**——``log_barrier_kcal``（harness 从沙箱内绑定日志推出的值，
cloudra/benchmark/accuracy.py:8	  ``barrier_evidence``）与 ``declared_barrier_kcal``（交付
cloudra/benchmark/accuracy.py:9	  ``mechanism_graph.json`` 里该边声明的 ``barrier_kcal_mol``）；
cloudra/benchmark/accuracy.py:10	- 不带参考值且有效可达（``accessible`` 且经由的 TS 节点也可达，见
cloudra/benchmark/accuracy.py:11	  ``GoldMechanismGraph.edge_effectively_accessible``）→ 该边记 1.0（可算但
cloudra/benchmark/accuracy.py:12	  标答未给参考值，记满分）；
cloudra/benchmark/accuracy.py:13	- 标答规定不可达（边自身或经由的 TS 节点 ``accessible=false``）→ 该边记 0.0
cloudra/benchmark/accuracy.py:14	  （无 accuracy 分；此结构记账按面值计，不被 ``max_penalty_fraction``
cloudra/benchmark/accuracy.py:15	  封顶救回——封顶只作用于带参考值边的高斯偏离/定义差异罚分）。
cloudra/benchmark/accuracy.py:16	
cloudra/benchmark/accuracy.py:17	带参考值的边缺日志推算、缺声明、或声明在**任何合法组合**下都得不到支撑且
cloudra/benchmark/accuracy.py:18	证据也对不上参考 → 该边 0 分并置 fail-loud 结构化标记
cloudra/benchmark/accuracy.py:19	（``AccuracyReport.fail_loud``），由 runner 据此短路 composite（final 只保留
cloudra/benchmark/accuracy.py:20	图分；绝不用最终消息里的数字兜底）。声明支撑得住（或虽支撑不住、但交付证据
cloudra/benchmark/accuracy.py:21	里另有与参考一致的组合——即 agent 选了另一个构象/参考系定义、而正确的计算
cloudra/benchmark/accuracy.py:22	已在它的日志里）→ 走正常数值路径或按「答案定义差异」固定小扣分
cloudra/benchmark/accuracy.py:23	（`_DEFINITION_MISMATCH_EDGE_SCORE`），**不触发 fail-loud**（用户口径：这是
cloudra/benchmark/accuracy.py:24	答案定义问题，不是科学错误）。一致后与 gold 边值比，用高斯容差评分；逐边
cloudra/benchmark/accuracy.py:25	分数按 gold 边 ``weight`` 加权平均（权重越大 = 该边越容易算准/越可核实）。
cloudra/benchmark/accuracy.py:26	带参考值边的罚分汇总后按 ``max_penalty_fraction`` 封顶；无参考边的
cloudra/benchmark/accuracy.py:27	1.0/0.0 结构记账不封顶（见上）。
cloudra/benchmark/accuracy.py:28	
cloudra/benchmark/accuracy.py:29	若 gold 机理图没有任何边 → accuracy 不评分（final_score=1.0），沿用旧 null
cloudra/benchmark/accuracy.py:30	行为。
cloudra/benchmark/accuracy.py:31	
\end{lstlisting}
\clearpage
\begin{lstlisting}[firstnumber=32,numbers=left,numberstyle=\tiny\ttfamily,numbersep=5pt]
cloudra/benchmark/accuracy.py:32	性质锚点（``GoldSpec.reference_properties``，键带单位）：读交付图物种节点的
cloudra/benchmark/accuracy.py:33	``properties`` 按键比对——数值型走高斯容差（σ 查模块级显式常量表
cloudra/benchmark/accuracy.py:34	``_PROPERTY_SIGMAS``，键名子串匹配，表外数值键 fail-loud 抛错逼显式登记）；
cloudra/benchmark/accuracy.py:35	字符串型（state_ordering 等）小写去空白后精确相等。多性质题取各键均值。
cloudra/benchmark/accuracy.py:36	
cloudra/benchmark/accuracy.py:37	hard 档定性选择性（``score_qualitative_selectivity``）：任务带
cloudra/benchmark/accuracy.py:38	``gold.selectivity.favored_channel`` 时替代 accuracy 槽位（composite 的
cloudra/benchmark/accuracy.py:39	accuracy 权重给选择性分）——交付边中恰好一条 ``favored: true`` 且其
cloudra/benchmark/accuracy.py:40	``channel`` 等于 gold favored_channel → 1.0，否则 0.0；barrier spread 数值
cloudra/benchmark/accuracy.py:41	不评分。
cloudra/benchmark/accuracy.py:42	"""
cloudra/benchmark/accuracy.py:43	
cloudra/benchmark/accuracy.py:44	from __future__ import annotations
cloudra/benchmark/accuracy.py:45	
cloudra/benchmark/accuracy.py:46	import math
cloudra/benchmark/accuracy.py:47	from dataclasses import dataclass, field
cloudra/benchmark/accuracy.py:48	from typing import TYPE_CHECKING, Optional
cloudra/benchmark/accuracy.py:49	
cloudra/benchmark/accuracy.py:50	from cloudra.benchmark.schema import (
cloudra/benchmark/accuracy.py:51	    AccuracyRubric,
cloudra/benchmark/accuracy.py:52	    GoldMechanismGraph,
cloudra/benchmark/accuracy.py:53	    GoldSpec,
cloudra/benchmark/accuracy.py:54	    SelectivityReference,
cloudra/benchmark/accuracy.py:55	)
cloudra/benchmark/accuracy.py:56	
cloudra/benchmark/accuracy.py:57	if TYPE_CHECKING:
cloudra/benchmark/accuracy.py:58	    from cloudra.benchmark.barrier_evidence import BarrierEvidence
cloudra/benchmark/accuracy.py:59	
cloudra/benchmark/accuracy.py:60	# 性质锚点的高斯容差表（键名子串匹配，首个命中生效；表外数值键 fail-loud）。
cloudra/benchmark/accuracy.py:61	# 新增性质键必须在此显式登记 σ——没有默认兜底值。
cloudra/benchmark/accuracy.py:62	_PROPERTY_SIGMAS = (
cloudra/benchmark/accuracy.py:63	    ("bond_length_angstrom", 0.02),
cloudra/benchmark/accuracy.py:64	    ("cm-1", 50.0),
cloudra/benchmark/accuracy.py:65	    ("kcal_mol", 2.0),
cloudra/benchmark/accuracy.py:66	)
cloudra/benchmark/accuracy.py:67	
cloudra/benchmark/accuracy.py:68	# 「答案定义差异」的逐边记分：证据里已经含有与参考一致的组合（agent 只是
cloudra/benchmark/accuracy.py:69	# 报/选了另一个构象或参考系定义），属定义层面的不一致而非科学错误——按此
cloudra/benchmark/accuracy.py:70	# 固定小扣分计入加权，绝不触发 fail-loud（用户口径）。
cloudra/benchmark/accuracy.py:71	_DEFINITION_MISMATCH_EDGE_SCORE = 0.9
cloudra/benchmark/accuracy.py:72	
cloudra/benchmark/accuracy.py:73	
cloudra/benchmark/accuracy.py:74	@dataclass
cloudra/benchmark/accuracy.py:75	class AccuracyReport:
cloudra/benchmark/accuracy.py:76	    enabled: bool = True
cloudra/benchmark/accuracy.py:77	    log_barrier_kcal: float | None = None
cloudra/benchmark/accuracy.py:78	    declared_barrier_kcal: float | None = None
cloudra/benchmark/accuracy.py:79	    rds_edge_id: str = ""
cloudra/benchmark/accuracy.py:80	    sigma_kcal: float = 3.0
cloudra/benchmark/accuracy.py:81	    gaussian_score: float = 0.0          # in [0, 1], 1 = exact match
cloudra/benchmark/accuracy.py:82	    capped_penalty: float = 0.0          # penalty actually applied after cap
cloudra/benchmark/accuracy.py:83	    final_score: float = 0.0             # 1 - capped_penalty（缺证据/不一致直接边 0）
\end{lstlisting}
\clearpage
\begin{lstlisting}[firstnumber=84,numbers=left,numberstyle=\tiny\ttfamily,numbersep=5pt]
cloudra/benchmark/accuracy.py:84	    # fail-loud：任一参考边的声明/推算缺失或不一致时置 True（本层 accuracy
cloudra/benchmark/accuracy.py:85	    # 记 0），runner 据此短路 composite（final 只保留图分）
cloudra/benchmark/accuracy.py:86	    fail_loud: bool = False
cloudra/benchmark/accuracy.py:87	    # 逐边能垒得分：edge_id -> 高斯分（空 = 非能垒模式或本层未评分）
cloudra/benchmark/accuracy.py:88	    edge_scores: dict[str, float] = field(default_factory=dict)
cloudra/benchmark/accuracy.py:89	    # 性质锚点模式：逐键得分（空 = 能垒模式或本层未评分）
cloudra/benchmark/accuracy.py:90	    property_scores: dict[str, float] = field(default_factory=dict)
cloudra/benchmark/accuracy.py:91	    notes: list[str] = field(default_factory=list)
cloudra/benchmark/accuracy.py:92	
cloudra/benchmark/accuracy.py:93	
cloudra/benchmark/accuracy.py:94	def gaussian_barrier_score(agent: float, gold: float, sigma: float) -> float:
cloudra/benchmark/accuracy.py:95	    """exp(-(agent-gold)^2 / (2 sigma^2)) in [0, 1]."""
cloudra/benchmark/accuracy.py:96	    if sigma <= 0:
cloudra/benchmark/accuracy.py:97	        return 1.0 if agent == gold else 0.0
cloudra/benchmark/accuracy.py:98	    d = agent - gold
cloudra/benchmark/accuracy.py:99	    return math.exp(-(d * d) / (2.0 * sigma * sigma))
cloudra/benchmark/accuracy.py:100	
cloudra/benchmark/accuracy.py:101	
cloudra/benchmark/accuracy.py:102	def _effective_sigma(rubric: AccuracyRubric, gold_barrier_kcal: float) -> float:
cloudra/benchmark/accuracy.py:103	    """逐边对 gold 比较的高斯宽度：小能垒放宽（rubric 小能垒口径），
cloudra/benchmark/accuracy.py:104	    大能垒用基准 sigma。只用于对 gold 的比较，不动内部一致性核对。"""
cloudra/benchmark/accuracy.py:105	    if gold_barrier_kcal <= rubric.small_barrier_threshold_kcal:
cloudra/benchmark/accuracy.py:106	        return max(rubric.sigma_kcal_mol, rubric.small_barrier_sigma_kcal)
cloudra/benchmark/accuracy.py:107	    return rubric.sigma_kcal_mol
cloudra/benchmark/accuracy.py:108	
cloudra/benchmark/accuracy.py:109	
cloudra/benchmark/accuracy.py:110	def _property_sigma(key: str) -> float:
cloudra/benchmark/accuracy.py:111	    """性质键的高斯 σ：显式常量表子串匹配；表外数值键 fail-loud（逼显式登记）。"""
cloudra/benchmark/accuracy.py:112	    for token, sigma in _PROPERTY_SIGMAS:
cloudra/benchmark/accuracy.py:113	        if token in key:
cloudra/benchmark/accuracy.py:114	            return sigma
cloudra/benchmark/accuracy.py:115	    raise ValueError(
cloudra/benchmark/accuracy.py:116	        f"no sigma registered for property key {key!r}; "
cloudra/benchmark/accuracy.py:117	        "extend _PROPERTY_SIGMAS explicitly"
cloudra/benchmark/accuracy.py:118	    )
cloudra/benchmark/accuracy.py:119	
cloudra/benchmark/accuracy.py:120	
cloudra/benchmark/accuracy.py:121	def _normalize_text(value: str) -> str:
cloudra/benchmark/accuracy.py:122	    return " ".join(value.lower().split())
cloudra/benchmark/accuracy.py:123	
cloudra/benchmark/accuracy.py:124	
cloudra/benchmark/accuracy.py:125	def _score_properties(
cloudra/benchmark/accuracy.py:126	    delivered_graph: Optional[GoldMechanismGraph],
cloudra/benchmark/accuracy.py:127	    gold_properties: dict[str, float | str],
cloudra/benchmark/accuracy.py:128	    rep: AccuracyReport,
cloudra/benchmark/accuracy.py:129	) -> AccuracyReport:
cloudra/benchmark/accuracy.py:130	    """性质锚点模式：交付图物种节点的 `properties` 逐键对 gold 比对取均值。"""
cloudra/benchmark/accuracy.py:131	    scores: dict[str, float] = {}
cloudra/benchmark/accuracy.py:132	    for key, gold_val in gold_properties.items():
cloudra/benchmark/accuracy.py:133	        if delivered_graph is None:
cloudra/benchmark/accuracy.py:134	            rep.notes.append(f"property {key}: no delivered graph; scored 0")
cloudra/benchmark/accuracy.py:135	            scores[key] = 0.0
\end{lstlisting}
\clearpage
\begin{lstlisting}[firstnumber=136,numbers=left,numberstyle=\tiny\ttfamily,numbersep=5pt]
cloudra/benchmark/accuracy.py:136	            continue
cloudra/benchmark/accuracy.py:137	        declared = [
cloudra/benchmark/accuracy.py:138	            (n.id, n.properties[key])
cloudra/benchmark/accuracy.py:139	            for n in delivered_graph.nodes
cloudra/benchmark/accuracy.py:140	            if key in n.properties
cloudra/benchmark/accuracy.py:141	        ]
cloudra/benchmark/accuracy.py:142	        if not declared:
cloudra/benchmark/accuracy.py:143	            rep.notes.append(
cloudra/benchmark/accuracy.py:144	                f"property {key}: not declared on any species node; scored 0"
cloudra/benchmark/accuracy.py:145	            )
cloudra/benchmark/accuracy.py:146	            scores[key] = 0.0
cloudra/benchmark/accuracy.py:147	            continue
cloudra/benchmark/accuracy.py:148	        if len(declared) > 1:
cloudra/benchmark/accuracy.py:149	            rep.notes.append(
cloudra/benchmark/accuracy.py:150	                f"property {key}: declared on multiple nodes "
cloudra/benchmark/accuracy.py:151	                f"{[nid for nid, _ in declared]}; using {declared[0][0]} "
cloudra/benchmark/accuracy.py:152	                "(first in graph order)"
cloudra/benchmark/accuracy.py:153	            )
cloudra/benchmark/accuracy.py:154	        node_id, val = declared[0]
cloudra/benchmark/accuracy.py:155	        if isinstance(gold_val, str):
cloudra/benchmark/accuracy.py:156	            if not isinstance(val, str):
cloudra/benchmark/accuracy.py:157	                rep.notes.append(
cloudra/benchmark/accuracy.py:158	                    f"property {key}: gold expects a string, node {node_id} "
cloudra/benchmark/accuracy.py:159	                    f"declared {type(val).__name__}; scored 0"
cloudra/benchmark/accuracy.py:160	                )
cloudra/benchmark/accuracy.py:161	                scores[key] = 0.0
cloudra/benchmark/accuracy.py:162	            else:
cloudra/benchmark/accuracy.py:163	                scores[key] = (
cloudra/benchmark/accuracy.py:164	                    1.0 if _normalize_text(val) == _normalize_text(gold_val) else 0.0
cloudra/benchmark/accuracy.py:165	                )
cloudra/benchmark/accuracy.py:166	                rep.notes.append(
cloudra/benchmark/accuracy.py:167	                    f"property {key}: node {node_id} declares {val!r} vs gold "
cloudra/benchmark/accuracy.py:168	                    f"{gold_val!r} -> {scores[key]}"
cloudra/benchmark/accuracy.py:169	                )
cloudra/benchmark/accuracy.py:170	        else:
cloudra/benchmark/accuracy.py:171	            if isinstance(val, bool) or not isinstance(val, (int, float)):
cloudra/benchmark/accuracy.py:172	                rep.notes.append(
cloudra/benchmark/accuracy.py:173	                    f"property {key}: gold expects a number, node {node_id} "
cloudra/benchmark/accuracy.py:174	                    f"declared {val!r}; scored 0"
cloudra/benchmark/accuracy.py:175	                )
cloudra/benchmark/accuracy.py:176	                scores[key] = 0.0
cloudra/benchmark/accuracy.py:177	            else:
cloudra/benchmark/accuracy.py:178	                sigma = _property_sigma(key)
cloudra/benchmark/accuracy.py:179	                scores[key] = gaussian_barrier_score(float(val), float(gold_val), sigma)
cloudra/benchmark/accuracy.py:180	                rep.notes.append(
cloudra/benchmark/accuracy.py:181	                    f"property {key}: node {node_id} declares {val} vs gold "
cloudra/benchmark/accuracy.py:182	                    f"{gold_val} (sigma {sigma}) -> {scores[key]:.4f}"
cloudra/benchmark/accuracy.py:183	                )
cloudra/benchmark/accuracy.py:184	    rep.property_scores = scores
cloudra/benchmark/accuracy.py:185	    rep.gaussian_score = sum(scores.values()) / len(scores)
cloudra/benchmark/accuracy.py:186	    rep.final_score = rep.gaussian_score
cloudra/benchmark/accuracy.py:187	    return rep
\end{lstlisting}
\clearpage
\begin{lstlisting}[firstnumber=188,numbers=left,numberstyle=\tiny\ttfamily,numbersep=5pt]
cloudra/benchmark/accuracy.py:188	
cloudra/benchmark/accuracy.py:189	
cloudra/benchmark/accuracy.py:190	def score_qualitative_selectivity(
cloudra/benchmark/accuracy.py:191	    delivered_graph: Optional[GoldMechanismGraph],
cloudra/benchmark/accuracy.py:192	    selectivity: SelectivityReference,
cloudra/benchmark/accuracy.py:193	) -> tuple[float, str]:
cloudra/benchmark/accuracy.py:194	    """hard 档定性选择性：交付边中恰好一条 favored=true 且 channel 命中 gold
cloudra/benchmark/accuracy.py:195	    favored_channel → 1.0；其余（无 favored 声明、多条 favored、channel 错）→ 0.0。
cloudra/benchmark/accuracy.py:196	    barrier spread 数值不评分。"""
cloudra/benchmark/accuracy.py:197	    if delivered_graph is None or not delivered_graph.edges:
cloudra/benchmark/accuracy.py:198	        return 0.0, (
cloudra/benchmark/accuracy.py:199	            "qualitative selectivity: no delivered graph edges; favored channel "
cloudra/benchmark/accuracy.py:200	            "not declared"
cloudra/benchmark/accuracy.py:201	        )
cloudra/benchmark/accuracy.py:202	    favored = [e for e in delivered_graph.edges if e.favored is True]
cloudra/benchmark/accuracy.py:203	    if not favored:
cloudra/benchmark/accuracy.py:204	        return 0.0, (
cloudra/benchmark/accuracy.py:205	            "qualitative selectivity: no edge declares `favored: true`; scored 0"
cloudra/benchmark/accuracy.py:206	        )
cloudra/benchmark/accuracy.py:207	    if len(favored) > 1:
cloudra/benchmark/accuracy.py:208	        return 0.0, (
cloudra/benchmark/accuracy.py:209	            "qualitative selectivity: multiple edges declare `favored: true` "
cloudra/benchmark/accuracy.py:210	            f"({[e.id for e in favored]}); the contract requires exactly one; scored 0"
cloudra/benchmark/accuracy.py:211	        )
cloudra/benchmark/accuracy.py:212	    edge = favored[0]
cloudra/benchmark/accuracy.py:213	    if edge.channel == selectivity.favored_channel:
cloudra/benchmark/accuracy.py:214	        return 1.0, (
cloudra/benchmark/accuracy.py:215	            f"qualitative selectivity: favored channel {edge.channel!r} matches "
cloudra/benchmark/accuracy.py:216	            "the reference; scored 1"
cloudra/benchmark/accuracy.py:217	        )
cloudra/benchmark/accuracy.py:218	    return 0.0, (
cloudra/benchmark/accuracy.py:219	        f"qualitative selectivity: favored channel {edge.channel!r} != reference "
cloudra/benchmark/accuracy.py:220	        f"{selectivity.favored_channel!r}; scored 0"
cloudra/benchmark/accuracy.py:221	    )
cloudra/benchmark/accuracy.py:222	
cloudra/benchmark/accuracy.py:223	
cloudra/benchmark/accuracy.py:224	def _find_matching_step(
cloudra/benchmark/accuracy.py:225	    gold_edge,
cloudra/benchmark/accuracy.py:226	    gold_nodes: dict,
cloudra/benchmark/accuracy.py:227	    steps: list,
cloudra/benchmark/accuracy.py:228	    delivered_graph: Optional[GoldMechanismGraph],
cloudra/benchmark/accuracy.py:229	    cfg,
cloudra/benchmark/accuracy.py:230	    match_threshold: float,
cloudra/benchmark/accuracy.py:231	):
cloudra/benchmark/accuracy.py:232	    """按**物种对应**找 gold 边的交付步（edge id 是题库内部约定，agent 不可知，
cloudra/benchmark/accuracy.py:233	    绝不参与匹配）：交付边的 from/to 端点物种与 gold 边**同向**端点分别匹配
cloudra/benchmark/accuracy.py:234	    （species_substitution_cost ≤ 阈值，与 Layer 1/graph GED 同口径；反向步
cloudra/benchmark/accuracy.py:235	    能垒口径不同不匹配）。多候选按 bond_changes 对称差最小、端点代价最小、
cloudra/benchmark/accuracy.py:236	    edge id 字典序取（确定性）。"""
cloudra/benchmark/accuracy.py:237	    from cloudra.benchmark.graph_similarity import species_substitution_cost
cloudra/benchmark/accuracy.py:238	
cloudra/benchmark/accuracy.py:239	    if delivered_graph is None:
\end{lstlisting}
\clearpage
\begin{lstlisting}[firstnumber=240,numbers=left,numberstyle=\tiny\ttfamily,numbersep=5pt]
cloudra/benchmark/accuracy.py:240	        return None
cloudra/benchmark/accuracy.py:241	    delivered_nodes = {n.id: n for n in delivered_graph.nodes}
cloudra/benchmark/accuracy.py:242	    delivered_edges = {e.id: e for e in delivered_graph.edges}
cloudra/benchmark/accuracy.py:243	    gold_from = gold_nodes.get(gold_edge.from_node)
cloudra/benchmark/accuracy.py:244	    gold_to = gold_nodes.get(gold_edge.to_node)
cloudra/benchmark/accuracy.py:245	    if gold_from is None or gold_to is None:
cloudra/benchmark/accuracy.py:246	        return None
cloudra/benchmark/accuracy.py:247	    gold_pairs = gold_edge.bond_changes.pairs()
cloudra/benchmark/accuracy.py:248	    candidates: list[tuple[int, float, str, object]] = []
cloudra/benchmark/accuracy.py:249	    for step in steps:
cloudra/benchmark/accuracy.py:250	        d_edge = delivered_edges.get(step.edge_id)
cloudra/benchmark/accuracy.py:251	        if d_edge is None:
cloudra/benchmark/accuracy.py:252	            continue
cloudra/benchmark/accuracy.py:253	        d_from = delivered_nodes.get(d_edge.from_node)
cloudra/benchmark/accuracy.py:254	        d_to = delivered_nodes.get(d_edge.to_node)
cloudra/benchmark/accuracy.py:255	        if d_from is None or d_to is None:
cloudra/benchmark/accuracy.py:256	            continue
cloudra/benchmark/accuracy.py:257	        cost_from = species_substitution_cost(gold_from.smiles, d_from.smiles, cfg)
cloudra/benchmark/accuracy.py:258	        cost_to = species_substitution_cost(gold_to.smiles, d_to.smiles, cfg)
cloudra/benchmark/accuracy.py:259	        if cost_from > match_threshold or cost_to > match_threshold:
cloudra/benchmark/accuracy.py:260	            continue
cloudra/benchmark/accuracy.py:261	        symdiff = len(gold_pairs ^ d_edge.bond_changes.pairs())
cloudra/benchmark/accuracy.py:262	        candidates.append((symdiff, cost_from + cost_to, step.edge_id, step))
cloudra/benchmark/accuracy.py:263	    if not candidates:
cloudra/benchmark/accuracy.py:264	        return None
cloudra/benchmark/accuracy.py:265	    candidates.sort(key=lambda item: (item[0], item[1], item[2]))
cloudra/benchmark/accuracy.py:266	    return candidates[0][3]
cloudra/benchmark/accuracy.py:267	
cloudra/benchmark/accuracy.py:268	
cloudra/benchmark/accuracy.py:269	def score_accuracy(
cloudra/benchmark/accuracy.py:270	    *,
cloudra/benchmark/accuracy.py:271	    evidence: "BarrierEvidence",
cloudra/benchmark/accuracy.py:272	    gold: GoldSpec,
cloudra/benchmark/accuracy.py:273	    rubric: AccuracyRubric | None = None,
cloudra/benchmark/accuracy.py:274	    delivered_graph: Optional[GoldMechanismGraph] = None,
cloudra/benchmark/accuracy.py:275	    cfg=None,
cloudra/benchmark/accuracy.py:276	    match_threshold: float = 0.2,
cloudra/benchmark/accuracy.py:277	) -> AccuracyReport:
cloudra/benchmark/accuracy.py:278	    """评分边集合 = 全部 gold 边：带参考值的边逐边核对（"声明得到日志支撑"），
cloudra/benchmark/accuracy.py:279	    再对 gold 比、按 gold 边 weight 加权平均；不带参考值的 accessible 边记 1.0、
cloudra/benchmark/accuracy.py:280	    不可达边记 0.0。gold.reference_properties 非空时走性质锚点模式（与能垒模式
cloudra/benchmark/accuracy.py:281	    互斥）。"""
cloudra/benchmark/accuracy.py:282	    rubric = rubric or AccuracyRubric()
cloudra/benchmark/accuracy.py:283	    rds = evidence.rds()
cloudra/benchmark/accuracy.py:284	    rep = AccuracyReport(
cloudra/benchmark/accuracy.py:285	        enabled=rubric.enabled,
cloudra/benchmark/accuracy.py:286	        sigma_kcal=rubric.sigma_kcal_mol,
cloudra/benchmark/accuracy.py:287	        log_barrier_kcal=rds.barrier_kcal_mol if rds is not None else None,
cloudra/benchmark/accuracy.py:288	        declared_barrier_kcal=rds.declared_kcal_mol if rds is not None else None,
cloudra/benchmark/accuracy.py:289	        rds_edge_id=rds.edge_id if rds is not None else "",
cloudra/benchmark/accuracy.py:290	    )
cloudra/benchmark/accuracy.py:291	
\end{lstlisting}
\clearpage
\begin{lstlisting}[firstnumber=292,numbers=left,numberstyle=\tiny\ttfamily,numbersep=5pt]
cloudra/benchmark/accuracy.py:292	    if not rubric.enabled:
cloudra/benchmark/accuracy.py:293	        rep.gaussian_score = 1.0
cloudra/benchmark/accuracy.py:294	        rep.final_score = 1.0
cloudra/benchmark/accuracy.py:295	        rep.notes.append("accuracy disabled for this task; scored 1.0 (computable but not graded)")
cloudra/benchmark/accuracy.py:296	        return rep
cloudra/benchmark/accuracy.py:297	    if gold.reference_properties:
cloudra/benchmark/accuracy.py:298	        return _score_properties(delivered_graph, gold.reference_properties, rep)
cloudra/benchmark/accuracy.py:299	    scored_edges = list(gold.mechanism_graph.edges)
cloudra/benchmark/accuracy.py:300	    if not scored_edges:
cloudra/benchmark/accuracy.py:301	        rep.gaussian_score = 1.0
cloudra/benchmark/accuracy.py:302	        rep.final_score = 1.0
cloudra/benchmark/accuracy.py:303	        rep.notes.append("gold mechanism graph carries no edges; accuracy not scored")
cloudra/benchmark/accuracy.py:304	        return rep
cloudra/benchmark/accuracy.py:305	
cloudra/benchmark/accuracy.py:306	    # --- 交付物双来源门禁：逐边核对（任一参考边缺失/不一致 -> fail-loud 标记） ---
cloudra/benchmark/accuracy.py:307	    from cloudra.benchmark.graph_similarity import GedConfig
cloudra/benchmark/accuracy.py:308	
cloudra/benchmark/accuracy.py:309	    cfg = cfg or GedConfig()
cloudra/benchmark/accuracy.py:310	    gold_nodes = {n.id: n for n in gold.mechanism_graph.nodes}
cloudra/benchmark/accuracy.py:311	    # 全零权重无法归一：退回逐边等权（避免除零；正常题库不出现）
cloudra/benchmark/accuracy.py:312	    equal_fallback = sum(e.weight for e in scored_edges) <= 0
cloudra/benchmark/accuracy.py:313	    if equal_fallback:
cloudra/benchmark/accuracy.py:314	        rep.notes.append(
cloudra/benchmark/accuracy.py:315	            "all scored-edge weights are zero; falling back to equal weighting"
cloudra/benchmark/accuracy.py:316	        )
cloudra/benchmark/accuracy.py:317	    weights = {e.id: (1.0 if equal_fallback else e.weight) for e in scored_edges}
cloudra/benchmark/accuracy.py:318	    weight_total = sum(weights.values())
cloudra/benchmark/accuracy.py:319	
cloudra/benchmark/accuracy.py:320	    failures: list[str] = []
cloudra/benchmark/accuracy.py:321	    weighted_sum = 0.0
cloudra/benchmark/accuracy.py:322	    graded_weight = 0.0        # 带参考值边的权重和（高斯/定义差异罚分只从这里出）
cloudra/benchmark/accuracy.py:323	    graded_score_sum = 0.0     # Σ w·g（带参考值边）
cloudra/benchmark/accuracy.py:324	    structural_penalty = 0.0   # Σ w·(1-s)（无参考边的结构记账：可算未给参考 s=1.0，标答不可达 s=0.0）
cloudra/benchmark/accuracy.py:325	    for edge in scored_edges:
cloudra/benchmark/accuracy.py:326	        if edge.barrier_kcal_mol is None:
cloudra/benchmark/accuracy.py:327	            if gold.mechanism_graph.edge_effectively_accessible(edge):
cloudra/benchmark/accuracy.py:328	                rep.edge_scores[edge.id] = 1.0
cloudra/benchmark/accuracy.py:329	                weighted_sum += weights[edge.id] * 1.0
cloudra/benchmark/accuracy.py:330	                rep.notes.append(
cloudra/benchmark/accuracy.py:331	                    f"edge {edge.id}: accessible but the gold reference carries no "
cloudra/benchmark/accuracy.py:332	                    "reference barrier; scored 1.0"
cloudra/benchmark/accuracy.py:333	                )
cloudra/benchmark/accuracy.py:334	            else:
cloudra/benchmark/accuracy.py:335	                rep.edge_scores[edge.id] = 0.0
cloudra/benchmark/accuracy.py:336	                structural_penalty += weights[edge.id]
cloudra/benchmark/accuracy.py:337	                rep.notes.append(
cloudra/benchmark/accuracy.py:338	                    f"edge {edge.id}: effectively inaccessible in the gold "
cloudra/benchmark/accuracy.py:339	                    "reference (the edge itself or its via transition state is "
cloudra/benchmark/accuracy.py:340	                    "declared accessible=false); no accuracy score (0.0)"
cloudra/benchmark/accuracy.py:341	                )
cloudra/benchmark/accuracy.py:342	            continue
cloudra/benchmark/accuracy.py:343	        step = _find_matching_step(
\end{lstlisting}
\clearpage
\begin{lstlisting}[firstnumber=344,numbers=left,numberstyle=\tiny\ttfamily,numbersep=5pt]
cloudra/benchmark/accuracy.py:344	            edge, gold_nodes, evidence.steps, delivered_graph, cfg, match_threshold
cloudra/benchmark/accuracy.py:345	        )
cloudra/benchmark/accuracy.py:346	        if step is None:
cloudra/benchmark/accuracy.py:347	            rep.edge_scores[edge.id] = 0.0
cloudra/benchmark/accuracy.py:348	            failures.append(
cloudra/benchmark/accuracy.py:349	                f"gold edge {edge.id} carries a reference barrier but no delivered "
cloudra/benchmark/accuracy.py:350	                "step matches its endpoint species"
cloudra/benchmark/accuracy.py:351	            )
cloudra/benchmark/accuracy.py:352	            continue
cloudra/benchmark/accuracy.py:353	        graded_weight += weights[edge.id]
cloudra/benchmark/accuracy.py:354	        if step.edge_id != edge.id:
cloudra/benchmark/accuracy.py:355	            rep.notes.append(
cloudra/benchmark/accuracy.py:356	                f"gold edge {edge.id} matched delivered step {step.edge_id} "
cloudra/benchmark/accuracy.py:357	                "by endpoint species"
cloudra/benchmark/accuracy.py:358	            )
cloudra/benchmark/accuracy.py:359	        if step.declared_kcal_mol is None:
cloudra/benchmark/accuracy.py:360	            rep.edge_scores[edge.id] = 0.0
cloudra/benchmark/accuracy.py:361	            failures.append(
cloudra/benchmark/accuracy.py:362	                "delivered mechanism_graph.json declares no barrier_kcal_mol on "
cloudra/benchmark/accuracy.py:363	                f"step edge {edge.id}"
cloudra/benchmark/accuracy.py:364	            )
cloudra/benchmark/accuracy.py:365	            continue
cloudra/benchmark/accuracy.py:366	        declared = step.declared_kcal_mol
cloudra/benchmark/accuracy.py:367	        # 该边全部合法组合（TS 日志 × 起点极小值日志；含策略选中的那组）。
cloudra/benchmark/accuracy.py:368	        # 声明支撑得住哪一组、哪一组能对上参考，决定这条边怎么算分：
cloudra/benchmark/accuracy.py:369	        # 支撑得住 → 正常数值路径；支撑不住但证据里有能对参考的 → 答案定义/
cloudra/benchmark/accuracy.py:370	        # 构象指认差异（小扣分，不是科学错误）；两者皆无 → 无据声明 fail-loud。
cloudra/benchmark/accuracy.py:371	        candidates = evidence.alternatives.get(step.edge_id) or [step]
cloudra/benchmark/accuracy.py:372	        sigma_eff = _effective_sigma(rubric, edge.barrier_kcal_mol)
cloudra/benchmark/accuracy.py:373	        supported = [
cloudra/benchmark/accuracy.py:374	            c for c in candidates
cloudra/benchmark/accuracy.py:375	            if abs(declared - c.barrier_kcal_mol) <= rubric.sigma_kcal_mol
cloudra/benchmark/accuracy.py:376	        ]
cloudra/benchmark/accuracy.py:377	        matches_ref = [
cloudra/benchmark/accuracy.py:378	            c for c in candidates
cloudra/benchmark/accuracy.py:379	            if abs(c.barrier_kcal_mol - edge.barrier_kcal_mol) <= sigma_eff
cloudra/benchmark/accuracy.py:380	        ]
cloudra/benchmark/accuracy.py:381	
cloudra/benchmark/accuracy.py:382	        if supported:
cloudra/benchmark/accuracy.py:383	            same_as_ref = any(
cloudra/benchmark/accuracy.py:384	                (c.ts_path, c.base_path) in {(m.ts_path, m.base_path) for m in matches_ref}
cloudra/benchmark/accuracy.py:385	                for c in supported
cloudra/benchmark/accuracy.py:386	            )
cloudra/benchmark/accuracy.py:387	            if matches_ref and not same_as_ref:
cloudra/benchmark/accuracy.py:388	                rep.edge_scores[edge.id] = _DEFINITION_MISMATCH_EDGE_SCORE
cloudra/benchmark/accuracy.py:389	                weighted_sum += weights[edge.id] * _DEFINITION_MISMATCH_EDGE_SCORE
cloudra/benchmark/accuracy.py:390	                graded_score_sum += weights[edge.id] * _DEFINITION_MISMATCH_EDGE_SCORE
cloudra/benchmark/accuracy.py:391	                ref = min(
cloudra/benchmark/accuracy.py:392	                    matches_ref,
cloudra/benchmark/accuracy.py:393	                    key=lambda c: (abs(c.barrier_kcal_mol - edge.barrier_kcal_mol),
cloudra/benchmark/accuracy.py:394	                                   c.ts_path, c.base_path),
cloudra/benchmark/accuracy.py:395	                )
\end{lstlisting}
\clearpage
\begin{lstlisting}[firstnumber=396,numbers=left,numberstyle=\tiny\ttfamily,numbersep=5pt]
cloudra/benchmark/accuracy.py:396	                rep.notes.append(
cloudra/benchmark/accuracy.py:397	                    f"edge {edge.id}: declared {declared:.3f} is supported by "
cloudra/benchmark/accuracy.py:398	                    f"{supported[0].ts_path} / {supported[0].base_path}, but the "
cloudra/benchmark/accuracy.py:399	                    f"delivered evidence also contains a combination consistent "
cloudra/benchmark/accuracy.py:400	                    f"with the reference ({ref.barrier_kcal_mol:.3f} vs gold "
cloudra/benchmark/accuracy.py:401	                    f"{edge.barrier_kcal_mol:.3f}, {ref.ts_path} / {ref.base_path}) "
cloudra/benchmark/accuracy.py:402	                    "-> answer-definition mismatch, scored "
cloudra/benchmark/accuracy.py:403	                    f"{_DEFINITION_MISMATCH_EDGE_SCORE}"
cloudra/benchmark/accuracy.py:404	                )
cloudra/benchmark/accuracy.py:405	                continue
cloudra/benchmark/accuracy.py:406	            chosen = step if step in supported else min(
cloudra/benchmark/accuracy.py:407	                supported,
cloudra/benchmark/accuracy.py:408	                key=lambda c: (abs(declared - c.barrier_kcal_mol), c.ts_path, c.base_path),
cloudra/benchmark/accuracy.py:409	            )
cloudra/benchmark/accuracy.py:410	            g = gaussian_barrier_score(chosen.barrier_kcal_mol, edge.barrier_kcal_mol, sigma_eff)
cloudra/benchmark/accuracy.py:411	            rep.edge_scores[edge.id] = g
cloudra/benchmark/accuracy.py:412	            weighted_sum += weights[edge.id] * g
cloudra/benchmark/accuracy.py:413	            graded_score_sum += weights[edge.id] * g
cloudra/benchmark/accuracy.py:414	            rep.notes.append(
cloudra/benchmark/accuracy.py:415	                f"edge {edge.id}: log-derived {chosen.barrier_kcal_mol:.3f} vs gold "
cloudra/benchmark/accuracy.py:416	                f"{edge.barrier_kcal_mol:.3f} (sigma {sigma_eff:.1f}) -> {g:.4f} "
cloudra/benchmark/accuracy.py:417	                f"(weight {weights[edge.id]})"
cloudra/benchmark/accuracy.py:418	            )
cloudra/benchmark/accuracy.py:419	            continue
cloudra/benchmark/accuracy.py:420	
cloudra/benchmark/accuracy.py:421	        if matches_ref:
cloudra/benchmark/accuracy.py:422	            ref = min(
cloudra/benchmark/accuracy.py:423	                matches_ref,
cloudra/benchmark/accuracy.py:424	                key=lambda c: (abs(c.barrier_kcal_mol - edge.barrier_kcal_mol),
cloudra/benchmark/accuracy.py:425	                               c.ts_path, c.base_path),
cloudra/benchmark/accuracy.py:426	            )
cloudra/benchmark/accuracy.py:427	            rep.edge_scores[edge.id] = _DEFINITION_MISMATCH_EDGE_SCORE
cloudra/benchmark/accuracy.py:428	            weighted_sum += weights[edge.id] * _DEFINITION_MISMATCH_EDGE_SCORE
cloudra/benchmark/accuracy.py:429	            graded_score_sum += weights[edge.id] * _DEFINITION_MISMATCH_EDGE_SCORE
cloudra/benchmark/accuracy.py:430	            rep.notes.append(
cloudra/benchmark/accuracy.py:431	                f"edge {edge.id}: declared {declared:.3f} is not supported by any "
cloudra/benchmark/accuracy.py:432	                "delivered log combination, but the evidence contains one "
cloudra/benchmark/accuracy.py:433	                f"consistent with the reference ({ref.barrier_kcal_mol:.3f} vs gold "
cloudra/benchmark/accuracy.py:434	                f"{edge.barrier_kcal_mol:.3f}, {ref.ts_path} / {ref.base_path}) "
cloudra/benchmark/accuracy.py:435	                "-> answer-definition mismatch, scored "
cloudra/benchmark/accuracy.py:436	                f"{_DEFINITION_MISMATCH_EDGE_SCORE}"
cloudra/benchmark/accuracy.py:437	            )
cloudra/benchmark/accuracy.py:438	            continue
cloudra/benchmark/accuracy.py:439	
cloudra/benchmark/accuracy.py:440	        rep.edge_scores[edge.id] = 0.0
cloudra/benchmark/accuracy.py:441	        failures.append(
cloudra/benchmark/accuracy.py:442	            "declared vs log-derived barrier disagreement on edge "
cloudra/benchmark/accuracy.py:443	            f"{edge.id} beyond sigma {rubric.sigma_kcal_mol}: declared "
cloudra/benchmark/accuracy.py:444	            f"{declared:.3f} vs logs {step.barrier_kcal_mol:.3f}"
cloudra/benchmark/accuracy.py:445	        )
cloudra/benchmark/accuracy.py:446	
cloudra/benchmark/accuracy.py:447	    if failures:
\end{lstlisting}
\clearpage
\begin{lstlisting}[firstnumber=448,numbers=left,numberstyle=\tiny\ttfamily,numbersep=5pt]
cloudra/benchmark/accuracy.py:448	        # fail-loud：任一参考边的声明/推算缺失或不一致 -> 置结构化标记、本层
cloudra/benchmark/accuracy.py:449	        # accuracy 记 0（不给部分分，也不套 max_penalty_fraction 封顶），
cloudra/benchmark/accuracy.py:450	        # composite 的短路由 runner 按 fail_loud 执行。
cloudra/benchmark/accuracy.py:451	        rep.notes.extend(f"{note}; scored 0" for note in failures)
cloudra/benchmark/accuracy.py:452	        rep.fail_loud = True
cloudra/benchmark/accuracy.py:453	        rep.gaussian_score = 0.0
cloudra/benchmark/accuracy.py:454	        rep.capped_penalty = 1.0
cloudra/benchmark/accuracy.py:455	        rep.final_score = 0.0
cloudra/benchmark/accuracy.py:456	        return rep
cloudra/benchmark/accuracy.py:457	
cloudra/benchmark/accuracy.py:458	    rep.gaussian_score = weighted_sum / weight_total
cloudra/benchmark/accuracy.py:459	    # max_penalty_fraction 只封顶带参考值边的罚分（高斯偏离/定义差异——兜合法的
cloudra/benchmark/accuracy.py:460	    # 方法学差异）；无参考边的结构记账（可算未给参考=1.0、标答不可达=0.0）是题面
cloudra/benchmark/accuracy.py:461	    # 口径的直接后果，按面值计、不参与封顶——否则 impossible 题的不可达边会被
cloudra/benchmark/accuracy.py:462	    # 封顶救回（0.7），设计上限（I 题满分 = 图 + 自洽）失守。
cloudra/benchmark/accuracy.py:463	    raw_penalty = 1.0 - rep.gaussian_score
cloudra/benchmark/accuracy.py:464	    graded_penalty = (graded_weight - graded_score_sum) / weight_total
cloudra/benchmark/accuracy.py:465	    capped_graded = min(graded_penalty, rubric.max_penalty_fraction)
cloudra/benchmark/accuracy.py:466	    rep.capped_penalty = capped_graded + structural_penalty / weight_total
cloudra/benchmark/accuracy.py:467	    rep.final_score = 1.0 - rep.capped_penalty
cloudra/benchmark/accuracy.py:468	    if graded_penalty > rubric.max_penalty_fraction:
cloudra/benchmark/accuracy.py:469	        rep.notes.append(
cloudra/benchmark/accuracy.py:470	            f"raw penalty {raw_penalty:.3f} capped at {rubric.max_penalty_fraction}"
cloudra/benchmark/accuracy.py:471	        )
cloudra/benchmark/accuracy.py:472	    return rep
cloudra/benchmark/api_pool.py:1	"""API endpoint pool for benchmark runs.
cloudra/benchmark/api_pool.py:2	
cloudra/benchmark/api_pool.py:3	Anti-429 for parallel benchmark execution: every LLM call picks a random
cloudra/benchmark/api_pool.py:4	endpoint from the pool (spreading load), and if the call still exhausts the
cloudra/benchmark/api_pool.py:5	client-level backoff retries with a rate-limit error, the next endpoint takes
cloudra/benchmark/api_pool.py:6	over (no repeats within one call).
cloudra/benchmark/api_pool.py:7	
cloudra/benchmark/api_pool.py:8	Scope by design: pool members share the *model* from the agent's settings —
cloudra/benchmark/api_pool.py:9	only the transport (base_url / api_key) is pinned per endpoint. Empty fields
cloudra/benchmark/api_pool.py:10	fall back to the agent's settings, so a pool can be just a list of extra keys
cloudra/benchmark/api_pool.py:11	for the same gateway. Non-429 errors are not retried here: they fail the turn
cloudra/benchmark/api_pool.py:12	(the runner's own turn-retry handles transient crashes as before).
cloudra/benchmark/api_pool.py:13	"""
cloudra/benchmark/api_pool.py:14	
cloudra/benchmark/api_pool.py:15	from __future__ import annotations
cloudra/benchmark/api_pool.py:16	
cloudra/benchmark/api_pool.py:17	import contextvars
cloudra/benchmark/api_pool.py:18	import json
cloudra/benchmark/api_pool.py:19	import logging
cloudra/benchmark/api_pool.py:20	import random
cloudra/benchmark/api_pool.py:21	from dataclasses import dataclass
cloudra/benchmark/api_pool.py:22	from pathlib import Path
cloudra/benchmark/api_pool.py:23	from typing import Any, Optional
cloudra/benchmark/api_pool.py:24	
cloudra/benchmark/api_pool.py:25	from cloudra.runtime_v2.adapters import LegacyLLMProvider
cloudra/benchmark/api_pool.py:26	
cloudra/benchmark/api_pool.py:27	logger = logging.getLogger(__name__)
\end{lstlisting}
\clearpage
\begin{lstlisting}[firstnumber=28,numbers=left,numberstyle=\tiny\ttfamily,numbersep=5pt]
cloudra/benchmark/api_pool.py:28	
cloudra/benchmark/api_pool.py:29	
cloudra/benchmark/api_pool.py:30	@dataclass(frozen=True)
cloudra/benchmark/api_pool.py:31	class Endpoint:
cloudra/benchmark/api_pool.py:32	    """One pool member; empty fields fall back to the agent's settings."""
cloudra/benchmark/api_pool.py:33	
cloudra/benchmark/api_pool.py:34	    base_url: str = ""
cloudra/benchmark/api_pool.py:35	    api_key: str = ""
cloudra/benchmark/api_pool.py:36	    label: str = ""  # 日志用；空则取 base_url
cloudra/benchmark/api_pool.py:37	
cloudra/benchmark/api_pool.py:38	    def display(self) -> str:
cloudra/benchmark/api_pool.py:39	        return self.label or self.base_url or "(settings default)"
cloudra/benchmark/api_pool.py:40	
cloudra/benchmark/api_pool.py:41	
cloudra/benchmark/api_pool.py:42	# 当前调用钉住的端点。contextvars 按 asyncio task / 线程隔离——并行跑题时
cloudra/benchmark/api_pool.py:43	# 共享同一个 provider 实例也不会串端点。
cloudra/benchmark/api_pool.py:44	_current_endpoint: contextvars.ContextVar[Optional[Endpoint]] = contextvars.ContextVar(
cloudra/benchmark/api_pool.py:45	    "benchmark_api_pool_endpoint", default=None
cloudra/benchmark/api_pool.py:46	)
cloudra/benchmark/api_pool.py:47	
cloudra/benchmark/api_pool.py:48	
cloudra/benchmark/api_pool.py:49	def load_endpoints(path: Path | str) -> list[Endpoint]:
cloudra/benchmark/api_pool.py:50	    """Load a pool JSON: [{"base_url": ..., "api_key": ..., "label": ...}, ...]."""
cloudra/benchmark/api_pool.py:51	    raw = json.loads(Path(path).read_text(encoding="utf-8"))
cloudra/benchmark/api_pool.py:52	    if not isinstance(raw, list) or not raw:
cloudra/benchmark/api_pool.py:53	        raise ValueError(f"api pool file must be a non-empty JSON list: {path}")
cloudra/benchmark/api_pool.py:54	    endpoints = []
cloudra/benchmark/api_pool.py:55	    for i, item in enumerate(raw):
cloudra/benchmark/api_pool.py:56	        if not isinstance(item, dict):
cloudra/benchmark/api_pool.py:57	            raise ValueError(f"api pool entry #{i} must be an object: {item!r}")
cloudra/benchmark/api_pool.py:58	        endpoints.append(
cloudra/benchmark/api_pool.py:59	            Endpoint(
cloudra/benchmark/api_pool.py:60	                base_url=str(item.get("base_url") or ""),
cloudra/benchmark/api_pool.py:61	                api_key=str(item.get("api_key") or ""),
cloudra/benchmark/api_pool.py:62	                label=str(item.get("label") or ""),
cloudra/benchmark/api_pool.py:63	            )
cloudra/benchmark/api_pool.py:64	        )
cloudra/benchmark/api_pool.py:65	    return endpoints
cloudra/benchmark/api_pool.py:66	
cloudra/benchmark/api_pool.py:67	
cloudra/benchmark/api_pool.py:68	def _is_rate_limit(exc: BaseException) -> bool:
cloudra/benchmark/api_pool.py:69	    try:
cloudra/benchmark/api_pool.py:70	        import openai
cloudra/benchmark/api_pool.py:71	    except Exception:  # pragma: no cover - openai 是核心依赖
cloudra/benchmark/api_pool.py:72	        return False
cloudra/benchmark/api_pool.py:73	    return isinstance(exc, openai.RateLimitError)
cloudra/benchmark/api_pool.py:74	
cloudra/benchmark/api_pool.py:75	
cloudra/benchmark/api_pool.py:76	class PooledLLMProvider(LegacyLLMProvider):
cloudra/benchmark/api_pool.py:77	    """LegacyLLMProvider whose transport endpoint is pinned per call from a pool."""
cloudra/benchmark/api_pool.py:78	
cloudra/benchmark/api_pool.py:79	    def __init__(self, endpoints: list[Endpoint], *, rng: Any = None) -> None:
\end{lstlisting}
\clearpage
\begin{lstlisting}[firstnumber=80,numbers=left,numberstyle=\tiny\ttfamily,numbersep=5pt]
cloudra/benchmark/api_pool.py:80	        super().__init__()
cloudra/benchmark/api_pool.py:81	        self._endpoints = list(endpoints)
cloudra/benchmark/api_pool.py:82	        # rng 可注入（测试用固定顺序）；生产默认 Random()
cloudra/benchmark/api_pool.py:83	        self._rng = rng if rng is not None else random.Random()
cloudra/benchmark/api_pool.py:84	
cloudra/benchmark/api_pool.py:85	    def _make_client(self, settings):
cloudra/benchmark/api_pool.py:86	        ep = _current_endpoint.get()
cloudra/benchmark/api_pool.py:87	        if ep is None:
cloudra/benchmark/api_pool.py:88	            return super()._make_client(settings)
cloudra/benchmark/api_pool.py:89	        from cloudra.infra.llm_client import LLMClient
cloudra/benchmark/api_pool.py:90	
cloudra/benchmark/api_pool.py:91	        return LLMClient(
cloudra/benchmark/api_pool.py:92	            api_key=ep.api_key or settings.api_key,
cloudra/benchmark/api_pool.py:93	            base_url=ep.base_url or settings.base_url,
cloudra/benchmark/api_pool.py:94	            timeout=settings.timeout,
cloudra/benchmark/api_pool.py:95	        )
cloudra/benchmark/api_pool.py:96	
cloudra/benchmark/api_pool.py:97	    async def complete(self, agent, messages, tools, on_delta=None):
cloudra/benchmark/api_pool.py:98	        if not self._endpoints:
cloudra/benchmark/api_pool.py:99	            return await super().complete(agent, messages, tools, on_delta)
cloudra/benchmark/api_pool.py:100	        # 每次调用随机洗牌端点顺序：既摊初始负载，又保证 429 时不重复试同一个
cloudra/benchmark/api_pool.py:101	        order = self._rng.sample(self._endpoints, k=len(self._endpoints))
cloudra/benchmark/api_pool.py:102	        last_exc: Optional[BaseException] = None
cloudra/benchmark/api_pool.py:103	        for ep in order:
cloudra/benchmark/api_pool.py:104	            token = _current_endpoint.set(ep)
cloudra/benchmark/api_pool.py:105	            try:
cloudra/benchmark/api_pool.py:106	                return await super().complete(agent, messages, tools, on_delta)
cloudra/benchmark/api_pool.py:107	            except BaseException as exc:
cloudra/benchmark/api_pool.py:108	                if not _is_rate_limit(exc):
cloudra/benchmark/api_pool.py:109	                    raise
cloudra/benchmark/api_pool.py:110	                last_exc = exc
cloudra/benchmark/api_pool.py:111	                logger.warning("endpoint %s rate-limited (429); trying next", ep.display())
cloudra/benchmark/api_pool.py:112	            finally:
cloudra/benchmark/api_pool.py:113	                _current_endpoint.reset(token)
cloudra/benchmark/api_pool.py:114	        assert last_exc is not None
cloudra/benchmark/api_pool.py:115	        raise last_exc
cloudra/benchmark/backend.py:1	"""runtime_v2-backed agent backend for the benchmark runner.
cloudra/benchmark/backend.py:2	
cloudra/benchmark/backend.py:3	Per task this spins up an in-process runtime_v2 kernel (RuntimeCore +
cloudra/benchmark/backend.py:4	TurnPipeline + AgentRunner), spawns the task's coordinator agent with the task
cloudra/benchmark/backend.py:5	query against the per-task sandbox as the project workspace, and waits for the
cloudra/benchmark/backend.py:6	task to genuinely finish (see ``watch.py`` for the wait_background /
cloudra/benchmark/backend.py:7	wait_review state machine, GH #30). The method whitelist is enforced by
cloudra/benchmark/backend.py:8	appending `MethodPolicyStage` to the pipeline; compute jobs submitted via
cloudra/benchmark/backend.py:9	`submit_compute_job` are claimed and executed by a per-task in-process
cloudra/benchmark/backend.py:10	`SchedulerDaemon` (production runs it as a separate `cloudra scheduler`
cloudra/benchmark/backend.py:11	process instead).
cloudra/benchmark/backend.py:12	
cloudra/benchmark/backend.py:13	Process-wide isolation is the caller's job: the benchmark CLI points the
cloudra/benchmark/backend.py:14	global config's `data_dir` at the benchmark data dir before any tool runs, so
cloudra/benchmark/backend.py:15	session stores, compute jobs and MEMORY.md never touch the live conversation
cloudra/benchmark/backend.py:16	data dir. LLM credentials still come from the normal global config slots
\end{lstlisting}
\clearpage
\begin{lstlisting}[firstnumber=17,numbers=left,numberstyle=\tiny\ttfamily,numbersep=5pt]
cloudra/benchmark/backend.py:17	(coordinator/follower), unchanged.
cloudra/benchmark/backend.py:18	"""
cloudra/benchmark/backend.py:19	
cloudra/benchmark/backend.py:20	from __future__ import annotations
cloudra/benchmark/backend.py:21	
cloudra/benchmark/backend.py:22	import asyncio
cloudra/benchmark/backend.py:23	import logging
cloudra/benchmark/backend.py:24	import os
cloudra/benchmark/backend.py:25	import threading
cloudra/benchmark/backend.py:26	from pathlib import Path
cloudra/benchmark/backend.py:27	from typing import Any
cloudra/benchmark/backend.py:28	
cloudra/benchmark/backend.py:29	from cloudra.benchmark.policy_stage import MethodPolicyStage
cloudra/benchmark/backend.py:30	from cloudra.benchmark.schema import TaskSpec
cloudra/benchmark/backend.py:31	from cloudra.benchmark.watch import RunWatcher
cloudra/benchmark/backend.py:32	
cloudra/benchmark/backend.py:33	logger = logging.getLogger(__name__)
cloudra/benchmark/backend.py:34	
cloudra/benchmark/backend.py:35	class RuntimeV2Backend:
cloudra/benchmark/backend.py:36	    """`AgentBackend` implementation: run each task on the real runtime_v2.
cloudra/benchmark/backend.py:37	
cloudra/benchmark/backend.py:38	    `llm` / `tools` are injectable so tests can run the
cloudra/benchmark/backend.py:39	    whole wiring hermetically; production defaults are the legacy bridges.
cloudra/benchmark/backend.py:40	    """
cloudra/benchmark/backend.py:41	
cloudra/benchmark/backend.py:42	    def __init__(
cloudra/benchmark/backend.py:43	        self,
cloudra/benchmark/backend.py:44	        *,
cloudra/benchmark/backend.py:45	        run_timeout_s: float = 7200.0,
cloudra/benchmark/backend.py:46	        stuck_grace_s: float = 300.0,
cloudra/benchmark/backend.py:47	        poll_interval_s: float = 2.0,
cloudra/benchmark/backend.py:48	        llm: Any = None,
cloudra/benchmark/backend.py:49	        tools: Any = None,
cloudra/benchmark/backend.py:50	        start_scheduler: bool = True,
cloudra/benchmark/backend.py:51	        store_path: Any = None,
cloudra/benchmark/backend.py:52	    ) -> None:
cloudra/benchmark/backend.py:53	        self.run_timeout_s = run_timeout_s
cloudra/benchmark/backend.py:54	        self.stuck_grace_s = stuck_grace_s
cloudra/benchmark/backend.py:55	        self.poll_interval_s = poll_interval_s
cloudra/benchmark/backend.py:56	        self.llm = llm
cloudra/benchmark/backend.py:57	        self.tools = tools
cloudra/benchmark/backend.py:58	        self.start_scheduler = start_scheduler
cloudra/benchmark/backend.py:59	        # 事件级持久化（可选但 CLI 默认开）：全战役共享一个 runtime-v2.db，
cloudra/benchmark/backend.py:60	        # 只 append 不 restore——benchmark core 是一次性的，恢复语义无意义；
cloudra/benchmark/backend.py:61	        # 崩溃后事件历史与 token 遥测（cloudra stats tokens）仍可查。
cloudra/benchmark/backend.py:62	        self.store_path = Path(store_path) if store_path is not None else None
cloudra/benchmark/backend.py:63	        # 每次任务的 RuntimeCore 按 task_id 留存，供测试与调试内省事件日志；
cloudra/benchmark/backend.py:64	        # 并行跑题时多线程并发写，用锁保护。
cloudra/benchmark/backend.py:65	        self._cores: dict[str, Any] = {}
cloudra/benchmark/backend.py:66	        self._cores_lock = threading.Lock()
cloudra/benchmark/backend.py:67	
cloudra/benchmark/backend.py:68	    @property
\end{lstlisting}
\clearpage
\begin{lstlisting}[firstnumber=69,numbers=left,numberstyle=\tiny\ttfamily,numbersep=5pt]
cloudra/benchmark/backend.py:69	    def last_core(self) -> Any:
cloudra/benchmark/backend.py:70	        """最近一次任务的 core（串行调试便利）；并行时用 core_by_task()。"""
cloudra/benchmark/backend.py:71	        with self._cores_lock:
cloudra/benchmark/backend.py:72	            return next(reversed(self._cores.values())) if self._cores else None
cloudra/benchmark/backend.py:73	
cloudra/benchmark/backend.py:74	    def core_by_task(self, task_id: str) -> Any:
cloudra/benchmark/backend.py:75	        with self._cores_lock:
cloudra/benchmark/backend.py:76	            return self._cores.get(task_id)
cloudra/benchmark/backend.py:77	
cloudra/benchmark/backend.py:78	    def __call__(self, task: TaskSpec, sandbox: Path) -> None:
cloudra/benchmark/backend.py:79	        asyncio.run(self._run_task(task, Path(sandbox)))
cloudra/benchmark/backend.py:80	
cloudra/benchmark/backend.py:81	    async def _run_task(self, task: TaskSpec, sandbox: Path) -> None:
cloudra/benchmark/backend.py:82	        from cloudra.runtime_v2.adapters import (
cloudra/benchmark/backend.py:83	            LegacyLLMProvider,
cloudra/benchmark/backend.py:84	            LegacyPromptProvider,
cloudra/benchmark/backend.py:85	            LegacyToolExecutor,
cloudra/benchmark/backend.py:86	            legacy_default_runtime_settings,
cloudra/benchmark/backend.py:87	        )
cloudra/benchmark/backend.py:88	        from cloudra.runtime_v2.agent_runner import AgentRunner
cloudra/benchmark/backend.py:89	        from cloudra.runtime_v2.commands import CommandBus
cloudra/benchmark/backend.py:90	        from cloudra.runtime_v2.compute_notify import ComputeNotifyBroker
cloudra/benchmark/backend.py:91	        from cloudra.runtime_v2.core import RuntimeCore
cloudra/benchmark/backend.py:92	        from cloudra.runtime_v2.policies import default_stages
cloudra/benchmark/backend.py:93	        from cloudra.runtime_v2.schemas import Command
cloudra/benchmark/backend.py:94	        from cloudra.runtime_v2.turn_pipeline import TurnPipeline
cloudra/benchmark/backend.py:95	
cloudra/benchmark/backend.py:96	        sandbox.mkdir(parents=True, exist_ok=True)
cloudra/benchmark/backend.py:97	        sandbox = sandbox.resolve()
cloudra/benchmark/backend.py:98	
cloudra/benchmark/backend.py:99	        core = RuntimeCore(strict_events=True)
cloudra/benchmark/backend.py:100	        with self._cores_lock:
cloudra/benchmark/backend.py:101	            self._cores[task.task_id] = core
cloudra/benchmark/backend.py:102	        if self.store_path is not None:
cloudra/benchmark/backend.py:103	            from cloudra.runtime_v2.persistence import RuntimeStore
cloudra/benchmark/backend.py:104	
cloudra/benchmark/backend.py:105	            # 与 services/api_v2/app.py 同款接法；多题共享一个库文件，append
cloudra/benchmark/backend.py:106	            # 经 to_thread 卸载，WAL + busy_timeout 由 RuntimeStore 自带。
cloudra/benchmark/backend.py:107	            core.store = RuntimeStore(self.store_path)
cloudra/benchmark/backend.py:108	
cloudra/benchmark/backend.py:109	            async def persist_event(envelope) -> None:
cloudra/benchmark/backend.py:110	                # llm.chunk deltas 是瞬时 hub 事件：无回放价值，落库会灌爆事件表。
cloudra/benchmark/backend.py:111	                if envelope.type == "llm.chunk":
cloudra/benchmark/backend.py:112	                    return
cloudra/benchmark/backend.py:113	                await asyncio.to_thread(core.store.append_event, envelope)
cloudra/benchmark/backend.py:114	
cloudra/benchmark/backend.py:115	            core.event_hub.subscribe(persist_event)
cloudra/benchmark/backend.py:116	        core.runtime_defaults = legacy_default_runtime_settings()
cloudra/benchmark/backend.py:117	        stage = MethodPolicyStage(task.method_policy, base_dir=sandbox)
cloudra/benchmark/backend.py:118	        stages = [*default_stages(core), stage]
cloudra/benchmark/backend.py:119	        executor = self.tools or LegacyToolExecutor(core)
cloudra/benchmark/backend.py:120	        from cloudra.runtime_v2.invocation import InvocationService
\end{lstlisting}
\clearpage
\begin{lstlisting}[firstnumber=121,numbers=left,numberstyle=\tiny\ttfamily,numbersep=5pt]
cloudra/benchmark/backend.py:121	
cloudra/benchmark/backend.py:122	        invocations = InvocationService(
cloudra/benchmark/backend.py:123	            event_writer=core.event_writer, stages=stages, tools=executor, agents=core.agents
cloudra/benchmark/backend.py:124	        )
cloudra/benchmark/backend.py:125	        pipeline = TurnPipeline(
cloudra/benchmark/backend.py:126	            event_log=core.event_log,
cloudra/benchmark/backend.py:127	            stages=stages,
cloudra/benchmark/backend.py:128	            llm=self.llm or LegacyLLMProvider(),
cloudra/benchmark/backend.py:129	            tools=executor,
cloudra/benchmark/backend.py:130	            prompt_provider=LegacyPromptProvider(),
cloudra/benchmark/backend.py:131	            event_hub=core.event_hub,
cloudra/benchmark/backend.py:132	            notices=core.notices,
cloudra/benchmark/backend.py:133	            invocations=invocations,
cloudra/benchmark/backend.py:134	        )
cloudra/benchmark/backend.py:135	        core.runner = AgentRunner(core, pipeline)
cloudra/benchmark/backend.py:136	        # CLI uses a private RuntimeCore, so the live Web service's notify
cloudra/benchmark/backend.py:137	        # broker cannot see its sessions.  Start one here to deliver terminal
cloudra/benchmark/backend.py:138	        # compute notices and wake parked child runs.
cloudra/benchmark/backend.py:139	        broker = ComputeNotifyBroker(core, poll_interval=self.poll_interval_s)
cloudra/benchmark/backend.py:140	        scheduler = None
cloudra/benchmark/backend.py:141	        agent = None
cloudra/benchmark/backend.py:142	        bus = CommandBus(core)
cloudra/benchmark/backend.py:143	        try:
cloudra/benchmark/backend.py:144	            await broker.start()
cloudra/benchmark/backend.py:145	            if self.start_scheduler:
cloudra/benchmark/backend.py:146	                from cloudra.scheduler.daemon import SchedulerDaemon
cloudra/benchmark/backend.py:147	
cloudra/benchmark/backend.py:148	                # worker_id 含 pid+task_id：并行跑题时同进程多个实例不能共享
cloudra/benchmark/backend.py:149	                # 默认的 scheduler:<pid>，否则租约归属互相串。
cloudra/benchmark/backend.py:150	                scheduler = SchedulerDaemon(worker_id=f"scheduler:{os.getpid()}:{task.task_id}")
cloudra/benchmark/backend.py:151	                await scheduler.start()
cloudra/benchmark/backend.py:152	            project = (
cloudra/benchmark/backend.py:153	                await bus.execute(
cloudra/benchmark/backend.py:154	                    Command(
cloudra/benchmark/backend.py:155	                        type="project.create",
cloudra/benchmark/backend.py:156	                        payload={"name": f"benchmark:{task.task_id}", "workspace_dir": str(sandbox)},
cloudra/benchmark/backend.py:157	                    )
cloudra/benchmark/backend.py:158	                )
cloudra/benchmark/backend.py:159	            ).data
cloudra/benchmark/backend.py:160	            agent = (
cloudra/benchmark/backend.py:161	                await bus.execute(
cloudra/benchmark/backend.py:162	                    Command(
cloudra/benchmark/backend.py:163	                        type="agent.spawn",
cloudra/benchmark/backend.py:164	                        target={"project_id": project.id},
cloudra/benchmark/backend.py:165	                        payload={
cloudra/benchmark/backend.py:166	                            "role_id": task.prompt.role,
cloudra/benchmark/backend.py:167	                            "query": task.prompt.query,
cloudra/benchmark/backend.py:168	                            "settings": {
cloudra/benchmark/backend.py:169	                                "orchestration_mode": task.prompt.orchestration_mode.value,
cloudra/benchmark/backend.py:170	                                # 无人值守：破坏性工具审批走 yolo 自动通道
cloudra/benchmark/backend.py:171	                                "yolo": True,
cloudra/benchmark/backend.py:172	                                # 无人值守（GH #31）：隐藏并 fail-closed 拦截
\end{lstlisting}
\clearpage
\end{document}
